% !TEX root = ../main.tex
% !TEX encoding = UTF-8 Unicode
% !TEX TS-program = pdflatexmk
% !TEX spellcheck = en-US
% !BIB program = biber

\chapter{Spin Ruijsenaars--Schneider models from Coulomb branches}\label{chap:spinRSCB}

\begin{chapterabstract}
	\todo{tba}
\end{chapterabstract}

\section{Rational spin Ruijsenaars--Schneider model from 3d $\mathcal{N}=4$ gauge theory}

\subsection{GKLO representation of the homological Coulomb branch}

We proceed along the lines of \cite{gerasimov:2005} to construct a $\gamma$-deformed GKLO representation of the homological Coulomb branch algebra $\mathfrak{C}_{N,\ell}$. To this end, we give the following

\begin{definition}\label{def:ratDiffOps}
    Let $\mathfrak{A}_{N,\ell}$ be the commutative $\C\llbracket \gamma \rrbracket$-algebra
    \begin{equation}
        \mathfrak{A}_{N,\ell} \coloneq \C\llbracket \gamma \rrbracket[q_i^\alpha,(P_i^\alpha)^{\pm 1}][(q_i^\alpha-q_j^\beta)^{-1}]/J,
    \end{equation}
    where the generators $q_i^\alpha,P_i^\alpha$ have indices $i=1,\dots,N$ and $\alpha \in \Z$, we localize at the elements $q_i^\alpha-q_j^\beta$ for $(i,\alpha) \neq (j,\beta)$, and $J$ is the ideal generated by the cyclic relations
    \begin{equation}
        q_i^{\alpha+\ell} = q_i^\alpha - \gamma, \qquad P_i^{\alpha+\ell} = P_i^\alpha.
    \end{equation}
    We then make $\mathfrak{A}_{N,\ell}$ into a Poisson $\C\llbracket \gamma \rrbracket$-algebra via the log-canonical Poisson bracket
    \begin{equation}
        \{ q_i^\alpha, P_j^\beta \} = \delta_{ij} \delta^{\alpha\beta} P_j^\beta.
    \end{equation}
\end{definition}

\begin{prop}
    There is an injective homomorphism $\psi\colon \mathfrak{C}_{N,\ell} \to \mathfrak{A}_{N,\ell}/\gamma \mathfrak{A}_{N,\ell}$ of Poisson algebras that sends
    \begin{equation}
        \begin{aligned}
            \psi\colon \quad 
            q_i^\alpha \mapsto q_i^\alpha, \qquad
            u_i^{\alpha\pm} \mapsto \pm (P_i^\alpha)^{\pm 1} \frac{\chi^{\alpha \pm 1}(q_i^\alpha)}{\prod_{j \neq i} (q_i^\alpha-q_j^\alpha)} + \gamma \mathfrak{A}_{N,\ell}.
        \end{aligned}
    \end{equation}
\end{prop}

\begin{remark}
    Because of the existence of $\psi$, we justify abusing notation and writing
    \begin{equation}
        u_i^{\alpha\pm} \coloneq \pm (P_i^\alpha)^{\pm 1} \frac{\chi^{\alpha\pm1}(q_i^\alpha)}{\prod_{j \neq i} (q_i^\alpha-q_j^\alpha)} \in \mathfrak{A}_{N,\ell}
    \end{equation}
    as a shorthand. Henceforth we will only be working with these \emph{$\gamma$-deformed} monopole operators inside $\mathfrak{A}_{N,\ell}$.
\end{remark}

\begin{proof}
    When $\ell > 2$, we compute the Poisson brackets inside $\mathfrak{A}_{N,\ell}$ to be
    \begin{equation*}
        \{ u_i^{\alpha\pm},u_j^{\beta\pm} \} = \pm(1-\delta_{ij}\delta^{\alpha\beta})\kappa^{\alpha\beta} u_i^{\alpha\pm} u_j^{\beta\pm}
        \begin{cases}
            \frac{1}{q_i^\ell-q_j^{\ell-1}}, & (\alpha,\beta) = (0,\ell-1), \\
            \frac{1}{q_i^{\ell-1}-q_j^\ell}, & (\alpha,\beta) = (\ell-1,0), \\
            \frac{1}{q_i^\alpha-q_j^\beta}, & \text{otherwise}.
        \end{cases}
    \end{equation*}
    For $\ell=2$, we have
    \begin{equation*}
        \{ u_i^{\alpha\pm},u_j^{\beta\pm} \} = \pm(1-\delta_{ij}\delta^{\alpha\beta})\kappa^{\alpha\beta} u_i^{\alpha\pm} u_j^{\beta\pm}
        \begin{cases}
            \frac{1}{2} \big( \frac{1}{q_i^0-q_j^1} + \frac{1}{q_i^2-q_j^1} \big), & (\alpha,\beta) = (0,1), \\
            \frac{1}{2} \big( \frac{1}{q_i^1-q_j^0} + \frac{1}{q_i^1-q_j^2} \big), & (\alpha,\beta) = (1,0), \\
            \frac{1}{q_i^\alpha-q_j^\beta}, & \text{otherwise},
        \end{cases}
    \end{equation*}
    and for $\ell=1$, we have
    \begin{equation*}
        \{ u_i^{0\pm}, u_j^{0\pm}  \} = \pm (1-\delta_{ij})u_i^{0\pm} u_j^{0\pm} \bigg( \frac{2}{q_i^0-q_j^0} - \frac{1}{q_i^0-q_j^1} - \frac{1}{q_i^1-q_j^0} \bigg).
    \end{equation*}
    Since $\frac{1}{q_i^\ell - q_j^{\ell-1}} \equiv \frac{1}{q_i^0 - q_j^{\ell-1}} \mod \gamma \mathfrak{A}_{N,\ell}$, the result follows.
\end{proof}

\subsection{Affine Yangian of $\mathfrak{gl}_\ell$}

With this in hand, we may define the generating series
\begin{equation}
    e^\alpha(z) \coloneq \sum_{i=1}^N \frac{u_i^{\alpha+}}{z-q_i^\alpha} \in \mathfrak{A}_{N,\ell}\llbracket z^{-1} \rrbracket, \qquad f^\alpha(z) \coloneq \sum_{i=1}^N \frac{u_i^{\alpha-}}{z-q_i^\alpha} \in \mathfrak{A}_{N,\ell}\llbracket z^{-1} \rrbracket,
\end{equation}
as well as
\begin{equation}
    h^\alpha(z) \coloneq \frac{\chi^{\alpha+1}(z) \chi^{\alpha-1}(z)}{\chi^\alpha(z)^2} \in \mathfrak{A}_{N,\ell}\llbracket z^{-1} \rrbracket,
\end{equation}
and check that they satisfy the relations of the classical limit of the affine Yangian:

\begin{prop}
    The generating series $e^\alpha(z),f^\alpha(z),\chi^\alpha(z)$ define a representation of the classical limit of the $N$-truncated affine Yangian of $\mathfrak{gl}_\ell$ in the sense that $\chi^\alpha(z)$ is a polynomial of degree $N$ and the relations
    \begin{align}
        \{ \chi^\alpha(z), \chi^\beta(w) \} &=0, \\
        \{ e^\alpha(z),f^\beta(w) \}
        &= -\delta^{\alpha\beta} \frac{h^\alpha(z)-h^\beta(w)}{z-w}, \\
        \{ \chi^\alpha(z),e^\beta(w) \}
        &= \delta^{\alpha\beta} \chi^\alpha(z) \frac{e^\beta(z)-e^\beta(w)}{z-w}, \\
        \{ \chi^\alpha(z),f^\beta(w) \}
        &= -\delta^{\alpha\beta} \chi^\alpha(z) \frac{f^\beta(z)-f^\beta(w)}{z-w}
    \end{align}
    are satisfied.
\end{prop}

\begin{remark}
    This representation is nothing but the classical limit of the GKLO representation of the affine Yangian of $\mathfrak{gl}_\ell$ \cite{gerasimov:2005}. We note that $e^\alpha(z),f^\alpha(z)$ for $\alpha = 1,\dots,\ell-1$ and $\chi^\alpha(z)$ for $\alpha = 0,\dots,\ell-1$ generate the (finite) Yangian of $\mathfrak{gl}_\ell$ with quantum determinant
    \begin{equation}
        \operatorname{qdet}(z) = \prod_{\alpha=0}^{\ell-1} \frac{\chi^{\alpha+1}(z)}{\chi^\alpha(z)} = \frac{\chi^0(z+\gamma)}{\chi^0(z)}.
    \end{equation}
    When $\gamma = 0$, it follows that $\operatorname{qdet}(z) = 1$, which reduces us to the Yangian of $\mathfrak{sl}_\ell$. In that sense, $\gamma$ is the charge under the center of $\mathfrak{gl}_\ell$.
\end{remark}

\begin{corollary}
    The zero modes
    \begin{equation}
        E^\alpha \coloneq \frac{1}{2\pi\I} \oint_\infty e^\alpha(z) dz = \sum_{i=1}^N u_i^{\alpha+}, \qquad F^\alpha \coloneq \frac{1}{2\pi\I} \oint_\infty f^\alpha(z) dz = \sum_{i=1}^N u_i^{\alpha-},
    \end{equation}
    and
    \begin{equation}
        H^\alpha \coloneq \frac{1}{2\pi\I} \oint_\infty h^\alpha(z) dz = \sum_{i=1}^N (2q_i^\alpha - q_i^{\alpha+1} - q_i^{\alpha-1}),
    \end{equation}
    define a representation of $\widehat{\mathfrak{sl}}_\ell$ in the sense that
    \begin{align}
        \{ H^\alpha,H^\beta \} &= 0,
        \\[4pt]
        \{ H^\alpha,E^\beta \} &= \kappa^{\alpha\beta} E^\beta,
        \\[4pt]
        \{ H^\alpha,F^\beta \} &= -\kappa^{\alpha\beta} F^\beta,
        \\[4pt]
        \{ E^\alpha,F^\beta \} &= \delta^{\alpha\beta} H^\alpha.
    \end{align}
\end{corollary}

\subsection{$L$-operator algebra}

Our goal is to exhibit $\mathfrak{A}_{N,\ell}$ as the Poisson algebra of the rational spin RS model. To make such a connection, it is useful to have an $L$-operator algebra at our disposal.

\begin{definition}
    Introduce the one-site $L$-operators
    \begin{equation}
        L_{ij}^{\alpha\pm} \coloneq \frac{u_j^{\alpha+1,\pm}}{q_j^{\alpha+1}-q_i^\alpha} \in \mathfrak{A}_{N,\ell},
    \end{equation}
    as well as the total $L$-operator
    \begin{equation}\label{eq:totLRat}
        L \coloneq L^{0+} \cdots L^{\ell-1,+}.
    \end{equation}
\end{definition}

\begin{remark}
    We note that $L^{\alpha+\ell,\pm} = L^{\alpha\pm}$ by the cyclic relations of $\mathfrak{A}_{N,\ell}$.
\end{remark}

\begin{prop}
    The one-site $L$-operators satisfy the Poisson brackets
    \begin{equation}
        \{ L_1^{\alpha\pm}, L_2^{\beta\pm} \} = \pm(\delta^{\alpha\beta} r^\alpha L_1^{\alpha\pm} L_2^{\beta\pm} - \delta^{\alpha\beta} L_1^{\alpha\pm} L_2^{\beta\pm} \underline r^{\alpha+1} + \delta^{\alpha+1,\beta} L_1^{\alpha\pm} \bar r_{21}^\beta L_2^{\beta\pm} - \delta^{\alpha,\beta+1} L_2^{\beta\pm} \bar r^\alpha L_1^{\alpha\pm}),
    \end{equation}
    where we have used the matrices from \cite{arutyunov:1996}:
    \begin{align}\label{eq:ratMat}
        r^\alpha &\coloneq \sum_{i \neq j} \frac{1}{q_i^\alpha-q_j^\alpha} (e_{ii}-e_{ij}) \otimes (e_{jj} - e_{ji}), \\
        \bar r^\alpha &\coloneq \sum_{i \neq j} \frac{1}{q_i^\alpha-q_j^\alpha} (e_{ii}-e_{ij}) \otimes e_{jj},
        \\[4pt]
        \underline r^\alpha &\coloneq \sum_{i \neq j} \frac{1}{q_i^\alpha-q_j^\alpha} (e_{ij} \otimes e_{ji} - e_{ii} \otimes e_{jj}).
    \end{align}    
\end{prop}

\begin{corollary}
    The total $L$-operator satisfies
    \begin{equation}
        \{ L_1, L_2 \} = r^0 L_1 L_2 - L_1 L_2 \underline r^0 + L_1 \bar r_{21}^0 L_2 - L_2 \bar r^0 L_1,
    \end{equation}
    which reproduces the Poisson bracket of the Lax matrix from \cite{arutyunov:1998}.
\end{corollary}

\begin{proof}
    This follows from the Poisson algbera of the one-site $L$-operators and the identity
    \begin{equation*}
        r^\alpha + \bar r_{21}^\alpha - \bar r^\alpha - \underline r^\alpha = 0.
        \vspace{-1.5\baselineskip}
    \end{equation*}
\end{proof}

\begin{corollary}
    The Hamiltonians $H[n] \coloneq \operatorname{Tr} L^n$ are mutually Poisson-commuting.
\end{corollary}

\begin{prop}
    The Hamiltonians $H[n]$ are central with respect to $\widehat{\mathfrak{sl}}_\ell$:
    \begin{equation}
        \{ H[n], E^\alpha \} = 0, \qquad \{ H[n], F^\alpha \} = 0, \qquad \{ H[n], H^\alpha \} = 0.
    \end{equation}
\end{prop}

\begin{proof}
    Let us consider the bracket $\{ H[n],E^\alpha \}$ as an example. From the $L$-operator algebra, we derive the bracket
    \begin{equation*}
        \begin{aligned}
            \{ L_1^{\alpha+},u_2^{\beta+1,+} \}
            ={}& {-\delta^{\alpha\beta}} L_1^{\alpha+} u_2^{\beta+1,+} \underline r^{\alpha+1} - \delta^{\alpha,\beta+1} u_2^{\beta+1,+} \bar r^\alpha L_1^{\alpha+} - \delta^{\alpha\beta} Z L_1^{\alpha+} L_2^{\beta+} + \delta^{\alpha+1,\beta} L_1^{\alpha+} Z L_2^{\beta+}.
        \end{aligned}
    \end{equation*}
    with $Z = \sum_{i=1}^N e_{ii} \otimes e_i^t$. Then
    \begin{equation*}
        \begin{aligned}
            \{ H[n],E^\alpha \}
            ={}& \sum_{\mu=0}^{n\ell-1} \operatorname{Tr}_1 L_1^{0,+} \cdots L_1^{\mu-1,+} \{ L_1^\mu, u_2^{\alpha,+} \} L_1^{\mu+1,+} \cdots L_1^{n\ell-1,+} e_2 \\
            ={}& -\sum_{\mu=1}^{n\ell-1} \delta^{\mu\alpha} \operatorname{Tr}_1 L_1^{0,+} \cdots L_1^{\alpha-1,+} u_2^{\alpha,+} (\underline r^{\alpha} + \bar r^\alpha) L_1^{\alpha,+} \cdots L_1^{n\ell-1,+} e_2 - \delta^{\ell\alpha} \operatorname{Tr}_1 L_1^n u_2^{\alpha,+} (\underline r^{0}+\bar r^0) e_2 \\
            & + \delta^{0,\alpha-1} \operatorname{Tr}_1 (L_1^n Z L_2^{\alpha-1,+}-Z  L_2^{\alpha-1,+} L_1^n) e_2 \\
            ={}& 0,
        \end{aligned}
    \end{equation*}
    where we have used $(\underline r^\alpha + \bar r^\alpha) e_2 = 0$.
\end{proof}

\subsection{Superintegrability}

We have already seen that the Hamiltonians $H[n]$ Poisson-commute with the generators $E^\alpha,F^\alpha,H^\alpha$, which define a representation of the loop algebra $\widehat{\mathfrak{sl}}_\ell$. It turns out that this representation factors through a representation of the loop algebra $L(\mathfrak{gl}_\ell)$ with generators $J^{\alpha\beta}[n] \in \mathfrak{A}_{N,\ell}$, which can be expressed in terms of the $L$-operators. To see this, let
\begin{equation}
    V_{\alpha,\beta}^\pm \coloneq u^{\alpha\pm} L^{\alpha\pm} \cdots L^{\beta-1,\pm} e
\end{equation}
with $u^{\alpha\pm} = (u_1^{\alpha\pm},\dots,u_N^{\alpha\pm})$ and $e = (1,\dots,1)^t$. Then we can define
\begin{equation}
    J^{\alpha\beta}[0] \coloneq
    \begin{cases}
        V_{\beta,\alpha-1}^-, & \alpha > \beta \\
        \sum_{i=1}^N (q_i^\alpha - q_i^{\alpha-1}), & \alpha = \beta \\
        V_{\alpha,\beta-1}^+, & \alpha < \beta
    \end{cases},
\end{equation}
as well as
\begin{equation}
    J^{\alpha\beta}[-n] \coloneq V_{\beta,\alpha+n\ell-1}^-, \qquad J^{\alpha\beta}[n] \coloneq V_{\alpha,\beta+n\ell-1}^+.
\end{equation}
for $n > 0$. The coefficients $J^{\alpha\beta}[n]$ assemble into an $\ell \times \ell$ matrix $J[n]$.

\begin{prop}
    The generators $J^{\alpha\beta}[n]$ define a representation of the loop algebra $L(\mathfrak{gl}_\ell)$:
    \begin{equation}
        \{ J^{\alpha\beta}[n],J^{\mu\nu}[m] \} = \delta^{\mu\beta} J^{\alpha\nu}[n+m] - \delta^{\alpha\nu} J^{\mu\beta}[n+m].
\end{equation}
\end{prop}

\begin{remark}
    It was already clear from \cite[\S 6.6]{bullimore:2015} that $J^{\alpha\beta}[0]$ satisfies the relations of $\mathfrak{gl}_\ell$.
\end{remark}

\begin{lemma}
    We have $\operatorname{tr} J[n] = -\gamma H[n]$, where $\operatorname{tr}$ is the operation of taking the trace of an $\ell \times \ell$ matrix.
\end{lemma}

\begin{remark}
    In particular, the generators $J^{\alpha\beta}[n]$ define a representation of $L(\mathfrak{sl}_\ell)$ when $\gamma = 0$, which again exhibits $\gamma$ as the charge under the center of $\mathfrak{gl}_\ell$.
\end{remark}

\begin{proof}
    This follows from a telescopic argument similar to lemma \ref{lemma:ratLOp}.
\end{proof}

From this, we see that the Hamiltonians $H[n]$ are in the center of the subalgebra of $\mathfrak{A}_{N,\ell}$ generated by $J^{\alpha\beta}[n]$ and are thus part of a larger Poisson-commutative subalgebra given by the $N\ell$ algebraically independent Hamiltonians
\begin{equation}
    J^{\alpha\alpha}[n], \qquad n=1,\dots,N, \quad \alpha=1,\dots,\ell.
\end{equation}
Since $\mathfrak{A}_{N,\ell}$ has $2N\ell$ algebraically independent generators as a $\C\llbracket \gamma \rrbracket$-algebra, we conclude that $\mathfrak{A}_{N,\ell}$ describes an integrable model in the sense of the Liouville theorem. In fact, $\mathfrak{A}_{N,\ell}$ describes a \emph{superintegrable} model, since the central Hamiltonians $H[n]$ do not just commute among each other, but also commute with the loop algebra $L(\mathfrak{gl}_\ell)$.

\subsection{Equations of motion}

In this section, we show that the superintegrable model defined by the Poisson algebra $\mathfrak{A}_{N,\ell}$ and the Hamiltonians $J^{\alpha\alpha}[n]$ is the spin Ruijsenaars--Schneider model introduced by Krichever and Zabrodin \cite{krichever:1995}. To show the identification, we introduce
\begin{equation}
    a^\alpha \coloneq L^{0+} \cdots L^{\alpha-2,+} e, \qquad c^\alpha \coloneq u^{\alpha+} L^{\alpha+} \cdots L^{\ell-1,+},
\end{equation}
for $\alpha=1,\dots,\ell$. Notice that $a^1 = e$, in other words, we have the constraints $a_i^1 = 1$ for $i=1,\dots,N$. This should be contrasted with \cite{arutyunov:1998,chalykh:2020,arutyunov:2019,fairon:2026}, where the spin vectors satisfy the alternative constraint $\sum_\rho a_i^\rho = 1$. However, they differ only by an overall rescaling of the spin variables.

\begin{lemma}\label{lemma:ratLOp}
    The total $L$-operator can be written as
    \begin{equation}
        L_{ij} = -\frac{\sum_{\rho=1}^\ell a_i^\rho c_j^\rho}{q_i^0-q_j^\ell}.
\end{equation}
\end{lemma}

\begin{proof}
Indeed,
\begin{equation*}
    \begin{aligned}
        \sum_{\rho=1}^\ell &a_i^\rho c_j^\rho
        = \sum_{\rho=1}^\ell \sum_{k,l=1}^N (L^0 \cdots L^{\rho-2})_{il} u_k^{\rho+} (L^\rho \cdots L^{\ell-1})_{kj} \\
        ={}& \sum_{\rho=1}^\ell \sum_{k,l=1}^N (L^0 \cdots L^{\rho-2})_{il} L_{lk}^{\rho-1} (q_k^\rho-q_l^{\rho-1}) (L^\rho \cdots L^{\ell-1})_{kj} \\
        ={}& \sum_{\rho=1}^\ell \sum_{k=1}^N (L^0 \cdots L^{\rho-1})_{ik} (L^\rho \cdots L^{\ell-1})_{kj} q_k^\rho - \sum_{\rho=1}^\ell \sum_{l=1}^N (L^0 \cdots L^{\rho-2})_{il} (L^{\rho-1} \cdots L^{\ell-1})_{lj} q_l^{\rho-1} \\
        ={}& L_{ij} (q_j^\ell-q_i^0),
    \end{aligned}
\end{equation*}
which yields the result.
\end{proof}

\begin{theorem}
    The time evolution under the Hamiltonian $H \coloneq \gamma H[1]$ reproduces the equations of motion \eqref{eq:eom} with the identification $x_i = q_i^0$ and the rational potential $V(z) = \frac{1}{z}-\frac{1}{z+\gamma}$.
\end{theorem}

\begin{proof}
    Using the $L$-operator algebra, we find
    \begin{align*}
        \dot x_i &= \{ H,q_i^0 \} = -\gamma L_{ii},
        \\[7pt]
        \ddot x_i
        &= \{ H, \{ H,q_i^0 \} \} = \gamma^2 \sum_{j(\neq i)} \frac{2}{q_i^0-q_j^0} L_{ij} L_{ji}, \\
        \dot a_i^\alpha
        &= \{ H,a_i^\alpha \} = \gamma \sum_{j(\neq i)} \frac{1}{q_i^0-q_j^0} (a_i^\alpha-a_j^\alpha) L_{ij}, \\
        \dot c_i^\alpha
        &= \{ H,c_i^\alpha \} = -\gamma \sum_{j(\neq i)} \frac{1}{q_i^0-q_j^0} (c_i^\alpha L_{ij} + c_j^\alpha L_{ji}).
    \end{align*}
    We then use lemma \ref{lemma:ratLOp} and the identities
    \begin{align*}
        \gamma \frac{1}{q_i^0-q_j^0} \frac{1}{q_i^0 - q_j^\ell} &= V(x_i-x_j), \\
        \gamma^2 \frac{2}{q_i^0-q_j^0} \frac{1}{q_i^0 - q_j^\ell} \frac{1}{q_j^0 - q_i^\ell} &= V(x_i-x_j)-V(x_j-x_i)
    \end{align*}
    to arrive at the desired equations of motion.
\end{proof}

\subsection{Poisson bracket of the spin variables}

\begin{prop}
    The Poisson brackets of the spins can be written as
    \begin{align}
        \{ q_i^0, a_j^\alpha \} ={}& 0, \qquad \{ q_i^0,c_j^\alpha \} = \delta_{ij} c_j^\alpha, \\[8pt]
        \{ a_i^\alpha,a_j^\beta \} ={}& \frac{1-\delta_{ij}}{q_i^0-q_j^0} (a_j^\alpha-a_i^\alpha) (a_i^\beta-a_j^\beta),
        \\[4pt]
        \{ a_i^\alpha,c_j^\beta \} ={}& \frac{1-\delta_{ij}}{q_i^0-q_j^0} (a_j^\alpha-a_i^\alpha) c_j^\beta - \delta^{1\beta} L_{ij} a_i^\alpha + \delta^{\alpha\beta} L_{ij},
        \\[4pt]
        \{ c_i^\alpha,c_j^\beta \} ={}& \frac{1-\delta_{ij}}{q_i^0-q_j^0} (c_i^\alpha c_j^\beta + c_j^\alpha c_i^\beta) - \delta^{1\alpha} L_{ji} c_j^\beta + \delta^{1\beta} L_{ij} c_i^\alpha.
    \end{align}
\end{prop}

\begin{proof}
    Let us consider the Poisson bracket $\{ a_1^\alpha,a_2^\beta \}$ as an example, where $1$ and $2$ label auxiliary spaces. Introduce the partial monodromies $L^{\alpha,\beta} \coloneq L^\alpha \cdots L^\beta$. Then
    \begin{equation*}
        \begin{aligned}
            \{ &a_1^\alpha,a_2^\beta \} = \sum_{\mu=0}^{\alpha-2} \sum_{\nu=0}^{\beta-2} L_1^{0,\mu-1} L_2^{0,\nu-1} \{ L_1^\mu, L_2^\nu \} L_1^{\mu+1,\alpha-2} L_2^{\nu+1,\beta-2} e_1 e_2 \\
            ={}& \sum_{\mu=0}^{\alpha-2} \sum_{\nu=0}^{\beta-2} \delta^{\mu\nu} L_1^{0,\mu-1} L_2^{0,\nu-1} r^\mu L_1^{\mu,\alpha-2} L_2^{\nu,\beta-2} e_1 e_2 - \sum_{\mu=1}^{\alpha-1} \sum_{\nu=1}^{\beta-1} \delta^{\mu\nu} L_1^{0,\mu-1} L_2^{0,\nu-1} \underline r^\mu L_1^{\mu,\alpha-2} L_2^{\nu,\beta-2} e_1 e_2 \\
            &+ \sum_{\mu=1}^{\alpha-1} \sum_{\nu=1}^{\beta-2} \delta^{\mu\nu} L_1^{0,\mu-1} L_2^{0,\nu-1} \bar r_{21}^\mu L_1^{\mu,\alpha-2} L_2^{\nu,\beta-2} e_1 e_2 - \sum_{\mu=1}^{\alpha-2} \sum_{\nu=1}^{\beta-1} \delta^{\mu\nu} L_1^{0,\mu-1} L_2^{0,\nu-1} \bar r^\mu L_1^{\mu,\alpha-2} L_2^{\nu,\beta-2} e_1 e_2 \\
            ={}& \delta_{\alpha>1,\beta>1} r^0 a_1^\alpha a_2^\beta = r^0 a_1^\alpha a_2^\beta,
        \end{aligned}
    \end{equation*}
    where we used the identities
    \begin{equation*}
        \underline r^\alpha e_1 e_2 = 0, \qquad (\underline r^\alpha + \bar r^\alpha) e_2 = 0, \qquad (\underline r^\alpha - \bar r_{21}^\alpha) e_1 = 0, \qquad r e_1 = 0, \qquad r e_2 = 0.
    \end{equation*}
    The other brackets follow similarly. In total, we arrive at
    \begin{align*}
        \{ a_1^\alpha,a_2^\beta \} ={}& r^0 a_1^\alpha a_2^\beta \\
        \{ a_1^\alpha,c_2^\beta \} ={}& {-c_2^\beta} \bar r^0 a_1^\alpha - \delta^{1\beta} Z a_1^\alpha L_2 + \delta^{\alpha\beta} Z e_1 L_2 \\
        \{ c_1^\alpha,c_2^\beta \} ={}& {-c_1^\alpha} c_2^\beta \underline r^0 + \delta^{1\beta} c_1^\alpha Z L_2 - \delta^{1\alpha} c_2^\beta Z_{21} L_1.
    \end{align*}
    where $Z = \sum_{i=1}^N e_{ii} \otimes e_i^t$. In components, these are the claimed Poisson brackets.
\end{proof}

\begin{remark}
    Due to the constraint $a_i^1 = 1$, the ``physical'' spin degrees of freedom of the rational spin RS model are actually restricted to $a^\alpha,c^\alpha$ for $\alpha > 1$. We note that they form a quadratic Poisson algebra.
\end{remark}

\begin{corollary}
The rescaling $\hat a_i^\alpha \coloneq a_i^\alpha/\sum_\rho a_i^\rho, \hat c_i^\alpha \coloneq c_i^\alpha \sum_\rho a_i^\rho, \hat L_{ij} \coloneq \frac{\sum_\rho \hat a_i^\rho \hat c_j^\rho}{q_i^0-q_j^0+\gamma}$ reproduces the Poisson brackets from \cite{arutyunov:1998}:
\begin{align}
    \{ q_i^0, \hat a_j^\alpha \} &= 0, \qquad \{ q_i^0, \hat c_j^\alpha \} = \delta_{ij} \hat c_j^\alpha, \\[8pt]
    \{ \hat a_i^\alpha, \hat a_j^\beta \} &= \frac{1-\delta_{ij}}{q_i^0-q_j^0} (\hat a_j^\alpha-\hat a_i^\alpha) (\hat a_i^\beta-\hat a_j^\beta),
    \\[4pt]
    \{ \hat a_i^\alpha, \hat c_j^\beta \} &= \frac{1-\delta_{ij}}{q_i^0-q_j^0} (\hat a_j^\alpha-\hat a_i^\alpha) \hat c_j^\beta + \hat a_i^\alpha \hat L_{ij} - \delta^{\alpha\beta} \hat L_{ij},
    \\[4pt]
    \{ \hat c_i^\alpha, \hat c_j^\beta \} &= \frac{1-\delta_{ij}}{q_i^0-q_j^0} (\hat c_i^\alpha \hat c_j^\beta + \hat c_j^\alpha \hat c_i^\beta) - \hat c_i^\alpha \hat L_{ij} + \hat c_j^\beta \hat L_{ji}.
\end{align}
\end{corollary}


\section{Trigonometric spin Ruijsenaars--Schneider model from 4d $\mathcal{N}=2$ gauge theory}

\subsection{GKLO representation of the $K$-theoretic Coulomb branch}

For the hyperbolic case, we proceed in essentially the same way, starting with a $t$-deformation of the GKLO presentation of the $K$-theoretic Coulomb branch written down in \cite{finkelberg:2019,tsymbaliuk:2023}. We will use essentially the same notation as in the rational case for this section, as no confusion should arise.

\begin{definition}\label{def:hypDiffOps}
    Let $\mathfrak{A}_{N,\ell}^K$ be the commutative $\C\llbracket t-1 \rrbracket$-algebra
    \begin{equation}
        \mathfrak{A}_{N,\ell}^K \coloneq \C\llbracket t-1 \rrbracket[(Q_i^\alpha)^{\pm 1/2},(P_i^\alpha)^{\pm 1}][((Q_i^\alpha/Q_j^\beta)^{1/2}-(Q_j^\beta/Q_i^\alpha)^{1/2})^{-1}]/J,
    \end{equation}
    where the generators $Q_i^\alpha,P_i^\alpha$ have indices $i=1,\dots,N$ and $\alpha \in \Z$, we localize at the elements $(Q_i^\alpha/Q_j^\beta)^{1/2}-(Q_j^\beta/Q_i^\alpha)^{1/2}$ for $(i,\alpha) \neq (j,\beta)$, and $J$ is the ideal generated by the cyclic relations
    \begin{equation}
        Q_i^{\alpha+\ell} = t Q_i^\alpha, \qquad P_i^{\alpha+\ell} = P_i^\alpha.
    \end{equation}
    We then make $\mathfrak{A}_{N,\ell}^K$ into a Poisson $\C\llbracket t-1 \rrbracket$-algebra via the doubly log-canonical Poisson bracket
    \begin{equation}
        \{ Q_i^\alpha,P_j^\beta \} = \delta_{ij} \delta^{\alpha\beta} Q_i^\alpha P_j^\beta.
    \end{equation}
\end{definition}

\begin{prop}
    There exists an injective homomorphism $\psi^K\colon \mathfrak{C}_{N,\ell}^K \to \mathfrak{A}_{N,\ell}^K/(t-1)\mathfrak{A}_{N,\ell}^K$ of Poisson algebras that sends
    \begin{equation}
        \psi^K\colon \quad Q_i^\alpha \mapsto Q_i^\alpha, \qquad u_i^{\alpha\pm} \mapsto \pm (P_i^\alpha)^{\pm 1} \frac{\chi^{\alpha \pm 1}(Q_i^\alpha)}{\prod_{j \neq i} ((Q_i^\alpha/Q_j^\alpha)^{1/2}-(Q_j^\alpha/Q_i^\alpha)^{1/2})} + (t-1) \mathfrak{A}_{N,\ell}^K.
    \end{equation}
\end{prop}

\begin{remark}
    We again use the existence of $\psi^K$ as a justification for writing
    \begin{equation}
        u_i^{\alpha\pm} \coloneq \pm (P_i^\alpha)^{\pm 1} \frac{\chi^{\alpha \pm 1}(Q_i^\alpha)}{\prod_{j \neq i} ((Q_i^\alpha/Q_j^\alpha)^{1/2}-(Q_j^\alpha/Q_i^\alpha)^{1/2})} \in \mathfrak{A}_{N,\ell}^K
    \end{equation}
    as a shorthand.
\end{remark}

\begin{proof}
    We compute the Poisson brackets inside $\mathfrak{A}_{N,\ell}^K$ for $\ell > 2$ to be
    \begin{equation*}
        \{ u_i^{\alpha\pm},u_j^{\beta\pm} \} = \pm(1-\delta_{ij}\delta^{\alpha\beta}) \kappa^{\alpha\beta} u_i^{\alpha\pm} u_j^{\beta\pm}
        \begin{cases}
            \frac{1}{2} \frac{Q_i^\ell+Q_j^{\ell-1}}{Q_i^\ell-Q_j^{\ell-1}}, & (\alpha,\beta) = (0,\ell-1) \\
            \frac{1}{2} \frac{Q_i^{\ell-1}+Q_j^\ell}{Q_i^{\ell-1}-Q_j^\ell}, & (\alpha,\beta) = (\ell-1,0) \\
            \frac{1}{2} \frac{Q_i^\alpha+Q_j^\beta}{Q_i^\alpha-Q_j^\beta}, & \text{otherwise},
        \end{cases}
    \end{equation*}
    while for $\ell=2$, we get
    \begin{equation*}
        \{ u_i^{\alpha\pm},u_j^{\beta\pm} \} = \pm(1-\delta_{ij}\delta^{\alpha\beta})\kappa^{\alpha\beta} u_i^{\alpha\pm} u_j^{\beta\pm}
        \begin{cases}
            \frac{1}{4} \frac{Q_i^0+Q_j^{\ell-1}}{Q_i^0-Q_j^{\ell-1}} + \frac{1}{4} \frac{Q_i^\ell+Q_j^{\ell-1}}{Q_i^\ell-Q_j^{\ell-1}}, & (\alpha,\beta) = (0,\ell-1) \\
            \frac{1}{4} \frac{Q_i^{\ell-1}+Q_j^0}{Q_i^{\ell-1}-Q_j^0} + \frac{1}{4} \frac{Q_i^{\ell-1}+Q_j^\ell}{Q_i^{\ell-1}-Q_j^\ell}, & (\alpha,\beta) = (\ell-1,0) \\
            \frac{1}{2} \frac{Q_i^\alpha+Q_j^\beta}{Q_i^\alpha-Q_j^\beta}, & \text{otherwise},
        \end{cases}
    \end{equation*}
    and for $\ell=1$, it is
    \begin{equation*}
        \{ u_i^{0\pm},u_j^{0\pm} \} = \pm(1-\delta_{ij})u_i^{0\pm} u_j^{0\pm} \bigg( \frac{Q_i^0+Q_j^0}{Q_i^0-Q_j^0} - \frac{1}{2} \frac{Q_i^1+Q_j^0}{Q_i^1-Q_j^0} - \frac{1}{2} \frac{Q_i^0+Q_j^1}{Q_i^0-Q_j^1} \bigg).
    \end{equation*}
    Since $\frac{Q_i^\ell+Q_j^{\ell-1}}{Q_i^\ell-Q_j^{\ell-1}} \equiv \frac{Q_i^0+Q_j^{\ell-1}}{Q_i^0-Q_j^{\ell-1}} \mod (t-1) \mathfrak{A}_{N,\ell}^K$, the result follows.
\end{proof}

\subsection{Quantum toroidal algebra of $\mathfrak{gl}_\ell$}

In the case of the $K$-theoretic Coulomb branch, the affine Yangian is replaced by the quantum toroidal algebra. To see this, we introduce the generating series
\begin{equation}
    e^\alpha(z) \coloneq \sum_{i=1}^N \frac{u_i^{\alpha+}}{1-Q_i^\alpha/z} \in \mathfrak{A}_{N,\ell}^K\llbracket z^{-1} \rrbracket, \qquad
    f^\alpha(z) \coloneq \sum_{i=1}^N \frac{u_i^{\alpha-}}{1-Q_i^\alpha/z} \in \mathfrak{A}_{N,\ell}^K\llbracket z^{-1} \rrbracket,
\end{equation}
as well as
\begin{equation}\label{eq:modKthGaugePoly}
    h^\alpha(z) \coloneq \frac{\tilde \chi^{\alpha+1}(z) \tilde \chi^{\alpha-1}(z)}{\tilde \chi^\alpha(z)^2} \prod_{i=1}^N \frac{Q_i^\alpha}{(Q_i^{\alpha+1} Q_i^{\alpha-1})^{1/2}}, \qquad \tilde \chi^\alpha(z) \coloneq \prod_{i=1}^N (1-Q_i^\alpha/z).
\end{equation}
We then expand them according to
\begin{equation}
    e^\alpha(z) = \sum_{r \geq 0} e_r^\alpha z^{-r}, \qquad f^\alpha(z) = \sum_{r \geq 0} f_r^\alpha z^{-r}, \qquad h^\alpha(z) = \sum_{r \geq 0} k_r^{\alpha+} z^{-r} = \sum_{r \geq 0} k_r^{\alpha-} z^r.
\end{equation}

\begin{prop}
    The generating series $e^\alpha(z),f^\alpha(z),\tilde \chi^\alpha(z),h^\alpha(z)$ define a representation of the classical limit of the $N$-truncated positive half of the quantum toroidal algebra of $\mathfrak{gl}_\ell$ in the sense that
    \begin{align}
        \{ k_r^\alpha, k_s^\beta \} &= 0,
        \\[8pt]
        \{ e_r^\alpha,f_s^\beta \}
        &= \delta^{\alpha\beta} (k_{r+s}^{\alpha+} - k_{r+s}^{\alpha-}),
        \\[2pt]
        \{ \tilde \chi^\alpha(z),e^\beta(w) \}
        &= \delta^{\alpha\beta} \tilde \chi^\alpha(z) \frac{e^\beta(z)-e^\beta(w)}{z/w-1},
        \\[2pt]
        \{ \tilde \chi^\alpha(z),f^\beta(w) \}
        &= -\delta^{\alpha\beta} \tilde \chi^\alpha(z) \frac{f^\beta(z)-f^\beta(w)}{z/w-1}.
    \end{align}
\end{prop}

\begin{remark}
    This is the classical limit of the GKLO representation of the quantum toroidal algebra of $\mathfrak{gl}_\ell$ discussed in \cite{tsymbaliuk:2023}.
\end{remark}

\begin{corollary}
    The zero modes
    \begin{equation}
        E^\alpha \coloneq e_0^\alpha = \sum_{i=1}^N u_i^{\alpha+}, \quad F^\alpha \coloneq f_0^\alpha = \sum_{i=1}^N u_i^{\alpha-}, \quad K^\alpha \coloneq k_0^{\alpha+} = (k_0^{\alpha-})^{-1} = \prod_{i=1}^N \frac{Q_i^\alpha}{(Q_i^{\alpha+1} Q_i^{\alpha-1})^{1/2}}
    \end{equation}
    define a representation of the classical limit of the quantum affine algebra in the sense that
    \begin{align}
        \{ K^\alpha,K^\beta \} &= 0,
        \\[4pt]
        \{ K^\alpha,E^\beta \} &= \tfrac{1}{2} \kappa^{\alpha\beta} K^\alpha E^\beta,
        \\[4pt]
        \{ K^\alpha, F^\beta\} &= -\tfrac{1}{2} \kappa^{\alpha\beta} K^\alpha F^\beta,
        \\[4pt]
        \{ E^\alpha,F^\beta \} &= \delta^{\alpha\beta} (K^\alpha - (K^\alpha)^{-1}).
    \end{align}
\end{corollary}

\subsection{$L$-operator algebra}

As in the rational case, we can write down an algebra of $L$-operators:

\begin{definition}
    Introduce the one-site $L$-operators
    \begin{equation}
        L_{ij}^\alpha \coloneq \frac{u_j^{\alpha+1,+}}{1-Q_j^{\alpha+1}/Q_i^\alpha} \in \mathfrak{A}_{N,\ell}^K
    \end{equation}
    as well as the total $L$-operator
    \begin{equation}\label{eq:totLHyp}
        L \coloneq L^0 \cdots L^{\ell-1}.
    \end{equation}
\end{definition}

\begin{prop}
The $L$-operators satisfy the Poisson bracket
\begin{equation}\label{eq:Lops}
    \{ L_1^\alpha, L_2^\beta \} = \delta^{\alpha\beta} r^\alpha L_1^\alpha L_2^\beta - \delta^{\alpha\beta} L_1^\alpha L_2^\beta \underline r^{\alpha+1} + \delta^{\alpha+1,\beta} L_1^\alpha \bar r_{21}^\beta L_2^\beta - \delta^{\alpha,\beta+1} L_2^\beta \bar r^\alpha L_1^\alpha,
\end{equation}
where we have used the matrices from \cite{arutyunov:2019b}, except that $\bar r^\alpha$ is shifted by $-\tfrac{1}{2}$:
\begin{align}\label{eq:hypMat}
    r^\alpha &\coloneq \sum_{i \neq j} \Big( \frac{1}{Q_i^\alpha/Q_j^\alpha-1} e_{ii} - \frac{1}{1-Q_j^\alpha/Q_i^\alpha} e_{ij} \Big) \otimes (e_{jj} - e_{ji}),
    \\
    \bar r^\alpha &\coloneq \sum_{i \neq j} \frac{1}{1-Q_j^\alpha/Q_i^\alpha} (e_{ii}-e_{ij}) \otimes e_{jj} - \frac{1}{2},
    \\[4pt]
    \underline r^\alpha &\coloneq \sum_{i \neq j} \frac{1}{1-Q_j^\alpha/Q_i^\alpha} (e_{ij} \otimes e_{ji} - e_{ii} \otimes e_{jj}).
\end{align}
\end{prop}

\begin{remark}
    The $r$-matrices $r^\alpha,\underline r^\alpha$ are not anti-symmetric, but satisfy
    \begin{equation}
        r^\alpha + r_{21}^\alpha = C - 1, \qquad \underline r^\alpha + \underline r_{21}^\alpha = C - 1.
    \end{equation}
    with $C = \sum_{ij} e_{ij} \otimes e_{ji}$.
\end{remark}

\begin{corollary}
    The total $L$-operator satisfies
    \begin{equation}
        \{ L_1, L_2 \} = r^0 L_1 L_2 - L_1 L_2 \underline r^0 + L_1 \bar r_{21}^0 L_2 - L_2 \bar r^0 L_1,
    \end{equation}
    which reproduces the Poisson bracket of the Lax matrix from \cite{arutyunov:2019b}.
\end{corollary}

\begin{proof}
    This follows from the Poisson bracket of the one-site $L$-operators using the identity
    \begin{equation*}
        r^\alpha + \bar r_{21}^\alpha - \bar r^\alpha - \underline r^\alpha = 0.
        \vspace{-1.5\baselineskip}
    \end{equation*}
\end{proof}

\begin{corollary}
    The Hamiltonians $H[n] \coloneq \operatorname{Tr} L^n$ are mutually Poisson commuting.
\end{corollary}

\begin{prop}
    The Hamiltonians $H[n]$ are central in the classical limit of the quantum affine algebra:
    \begin{equation}
        \{ H[n], E^\alpha \} = 0, \qquad \{ H[n], F^\alpha \} = 0, \qquad \{ H[n], K^\alpha \} = 0.
    \end{equation}
\end{prop}

\begin{proof}
    Let us consider the bracket $\{ H[n],E^\alpha \}$ as an example. From the $L$-operator algebra, we derive the bracket
    \begin{equation*}
        \begin{aligned}
            \{ L_1^\alpha,u_2^{\beta+1,+} \}
            ={}& {-\delta^{\alpha\beta}} L_1^\alpha u_2^{\beta+1,+} \underline r^{\alpha+1} - \delta^{\alpha,\beta+1} u_2^{\beta+1,+} (\bar r^\alpha+\tfrac{1}{2}) L_1^\alpha + \delta^{\alpha\beta} Z L_1^\alpha \tilde L_2^\beta - \delta^{\alpha+1,\beta} L_1^\alpha Z \tilde L_2^\beta.
        \end{aligned}
    \end{equation*}
    with $Z = \sum_{i=1}^N e_{ii} \otimes e_i^t$. Then
    \begin{equation*}
        \begin{aligned}
            \{ H[n],E^\alpha \}
            ={}& \sum_{\mu=0}^{n\ell-1} \operatorname{Tr}_1 L_1^0 \cdots L_1^{\mu-1} \{ L_1^\mu, u_2^{\alpha,+} \} L_1^{\mu+1} \cdots L_1^{n\ell-1} e_2 \\
            ={}& -\sum_{\mu=1}^{n\ell-1} \delta^{\mu\alpha} \operatorname{Tr}_1 L_1^0 \cdots L_1^{\alpha-1} u_2^{\alpha,+} (\underline r^{\alpha} + \bar r^\alpha) L_1^{\alpha} \cdots L_1^{n\ell-1} e_2 - \delta^{\ell\alpha} \operatorname{Tr}_1 L_1^n u_2^{\alpha,+} (\underline r^{0}+\bar r^0) e_2 \\
            &- \sum_{\mu=2}^{n\ell+1} \tfrac{1}{2} \delta^{\mu\alpha} \operatorname{Tr}_1 L_1^n u_2^{\alpha,+} e_2 + \delta^{0,\alpha-1} \operatorname{Tr}_1 (Z \tilde L_2^{\alpha-1} L_1^n-L_1^n Z \tilde L_2^{\alpha-1}) e_2 \\
            ={}& 0,
        \end{aligned}
    \end{equation*}
    where we have used $(\underline r^\alpha + \bar r^\alpha) e_2 = -\tfrac{1}{2} e_2$.
\end{proof}

In analogy with the rational case, this suggests that the central Hamiltonians $H[n]$ together with the rest of the Bethe subalgebra of the classical limit of the quantum affine algebra define a superintegrable system with the $K$-theoretic Coulomb branch as its phase space. In the next section, we identify this superintegrable model with the hyperbolic spin RS model.

\subsection{Equations of motion}

Let us study the equations of motion under the first Hamiltonian $H[1]$. To this end, we let $Q^\alpha \coloneq \operatorname{diag}(Q_1^\alpha,\dots,Q_N^\alpha)$ and $\tilde L^\alpha \coloneq (Q^\alpha)^{-1} L^\alpha Q^{\alpha+1}$. Then we define the spin vectors
\begin{equation}
    a^\alpha \coloneq L^0 \cdots L^{\alpha-2} e, \qquad c^\alpha \coloneq u^{\alpha+} \tilde L^\alpha \cdots \tilde L^{\ell-1}.
\end{equation}
for $\alpha=1,\dots,\ell$. Notice again that $a_i^1 = 1$ for $i=1,\dots,N$.

\begin{lemma}\label{lemma:hypLOp}
    The total $L$-operator can be written as
    \begin{equation}
    L_{ij} = \frac{\sum_{\rho=1}^\ell a_i^\rho c_j^\rho}{1-Q_j^\ell/Q_i^0}.
\end{equation}
\end{lemma}

\begin{proof}
    Indeed,
    \begin{equation*}
        \begin{aligned}
            \sum_{\rho=1}^\ell &a_i^\rho c_j^\rho
            = \sum_{\rho=1}^\ell \sum_{k,l=1}^N (L^0 \cdots L^{\rho-2})_{il} u_k^{\rho+}  Q_j^\ell/Q_k^\rho (L^\rho \cdots L^{\ell-1})_{kj} \\
            ={}& \sum_{\rho=1}^\ell \sum_{k,l=1}^N (L^0 \cdots L^{\rho-2})_{il} L_{lk}^{\rho-1} (Q_j^\ell/Q_k^\rho-Q_j^\ell/Q_l^{\rho-1}) (L^\rho \cdots L^{\ell-1})_{kj} \\
            ={}& \sum_{\rho=1}^\ell \sum_{k=1}^N (L^0 \cdots L^{\rho-1})_{ik} (L^\rho \cdots L^{\ell-1})_{kj} Q_j^\ell/Q_k^\rho - \sum_{\rho=1}^\ell \sum_{l=1}^N (L^0 \cdots L^{\rho-2})_{il} (L^{\rho-1} \cdots L^{\ell-1})_{lj} Q_j^\ell/Q_l^{\rho-1} \\
            ={}& L_{ij} (1 - Q_j^\ell/Q_i^0),
        \end{aligned}
    \end{equation*}
    which yields the result.
\end{proof}

\begin{theorem}
The time evolution under the Hamiltonian $H \coloneq (t-1) H[1]$ reproduces the equations of motion \eqref{eq:eom} with the identification $x_i = \log Q_i^0, \gamma = -{\log t}$ and the hyperbolic potential $V(z) = \tfrac{1}{2} \coth \tfrac{z}{2}-\tfrac{1}{2} \coth \tfrac{z+\gamma}{2}$.
\end{theorem}

\begin{proof}
    We use the $L$-operator algebra to derive the equations of motion
    \begin{align*}
        \dot x_i &= \{ H,Q_i^0 \} / Q_i^0 = -(t-1)L_{ii},
        \\[4pt]
        \ddot x_i &= \{ H, \{ H,Q_i^0 \} / Q_i^0 \} = (t-1)^2 \sum_{i \neq j} L_{ij} L_{ji} \frac{Q_i^0+Q_j^0}{Q_i^0-Q_j^0}, \\
        \dot a_i^\alpha &= \{ H,a_i^\alpha \} = -(t-1)\sum_{i \neq j} (a_i^\alpha-a_j^\alpha) \frac{L_{ij}}{1-Q_i^0/Q_j^0}, \\
        \dot c_i^\alpha &= \{ H,c_i^\alpha \} = (t-1)\sum_{i \neq j} \Big( c_i^\alpha \frac{L_{ij}}{1-Q_i^0/Q_j^0} - c_j^\alpha \frac{L_{ji}}{1-Q_j^0/Q_i^0} \Big).
    \end{align*}
    We then use lemma \ref{lemma:hypLOp} and the identities
    \begin{align*}
        (t-1)\frac{1}{1-Q_i^0/Q_j^0} \frac{1}{1-Q_j^\ell/Q_i^0} &= V(x_i-x_j), \\
        (t-1)^2\frac{Q_i^0+Q_j^0}{Q_i^0-Q_j^0} \frac{1}{1-Q_j^\ell/Q_i^0} \frac{1}{1-Q_i^\ell/Q_j^0} &= V(x_i-x_j)-V(x_j-x_i)
    \end{align*}
    to arrive at the desired equations of motion \eqref{eq:eom}.
\end{proof}

\subsection{Poisson bracket of the spin variables}

\begin{prop}
    The Poisson brackets of the spins can be written as
    \begin{align}
        \{ Q_i^0, a_j^\alpha \} ={}& 0, \qquad \{ Q_i^0, c_j^\alpha \} = \delta_{ij} Q_i^0 c_j^\alpha, \\[4pt]
        \{ a_i^\alpha,a_j^\beta \} ={}& (1-\delta_{ij}) \bigg( \frac{a_i^\alpha (a_j^\beta-a_i^\beta)}{Q_i^0/Q_j^0-1} - \frac{a_j^\alpha (a_j^\beta-a_i^\beta)}{1-Q_j^0/Q_i^0} \bigg) - \delta^{\alpha < \beta} a_j^\alpha a_i^\beta \\
        &+ \tfrac{1}{2} (2\delta^{1=\alpha<\beta} + \delta^{\alpha>\beta>1} + \delta^{1 < \alpha < \beta}) a_i^\alpha a_j^\beta, \nonumber
        \\[4pt]
        \{ a_i^\alpha,c_j^\beta \} ={}& (1-\delta_{ij}) \frac{(a_j^\alpha-a_i^\alpha) c_j^\beta}{1-Q_j^0/Q_i^0} - \delta^{\alpha\beta} \sum_{\rho=1}^{\alpha-1} a_i^\rho c_j^\rho + \delta^{1\beta} \tilde L_{ij} a_i^\alpha - \delta^{\alpha\beta} \tilde L_{ij} \\
        &+ \tfrac{1}{2}(1 - \delta^{1=\alpha < \beta} + \delta^{\alpha > \beta=1}-\delta^{\alpha\beta}) a_i^\alpha c_j^\beta, \nonumber
        \\[4pt]
        \{ c_i^\alpha,c_j^\beta \} ={}& (1-\delta_{ij}) \bigg( \frac{c_j^\alpha c_i^\beta}{Q_i^0/Q_j^0-1} + \frac{c_i^\alpha c_j^\beta}{1-Q_j^0/Q_i^0} \bigg) + \delta^{\alpha < \beta} c_j^\alpha c_i^\beta - \delta^{1\beta} c_i^\alpha \tilde L_{ij} + \delta^{1\alpha} c_j^\beta \tilde L_{ji} \\
        &- \tfrac{1}{2} (2 \delta^{\alpha > \beta=1} + \delta^{\alpha > \beta > 1} + \delta^{1<\alpha<\beta}) c_i^\alpha c_j^\beta, \nonumber
    \end{align}
    where
    \begin{equation}
        \tilde L_{ij} = \frac{\sum_{\rho=1}^\ell a_i^\rho c_j^\rho}{Q_i^0/Q_j^\ell-1}.
    \end{equation}
\end{prop}

\begin{proof}
    We again consider the bracket $\{ a_1^\alpha,a_2^\beta \}$ as an example and make use of the partial monodromies $L^{\alpha,\beta} \coloneq L^\alpha \cdots L^\beta$. Here $1$ and $2$ again label auxiliary spaces. Then
    \begin{equation*}
        \begin{aligned}
            \{ &a_1^\alpha,a_2^\beta \} = \sum_{\mu=0}^{\alpha-2} \sum_{\nu=0}^{\beta-2} L_1^{0,\mu-1} L_2^{0,\nu-1} \{ L_1^\mu, L_2^\nu \} L_1^{\mu+1,\alpha-2} L_2^{\nu+1,\beta-2} e_1 e_2 \\
            ={}& \sum_{\mu=0}^{\alpha-2} \sum_{\nu=0}^{\beta-2} \delta^{\mu\nu} L_1^{0,\mu-1} L_2^{0,\nu-1} r^\mu L_1^{\mu,\alpha-2} L_2^{\nu,\beta-2} e_1 e_2 - \sum_{\mu=1}^{\alpha-1} \sum_{\nu=1}^{\beta-1} \delta^{\mu\nu} L_1^{0,\mu-1} L_2^{0,\nu-1} \underline r^\mu L_1^{\mu,\alpha-2} L_2^{\nu,\beta-2} e_1 e_2 \\
            &+ \sum_{\mu=1}^{\alpha-1} \sum_{\nu=1}^{\beta-2} \delta^{\mu\nu} L_1^{0,\mu-1} L_2^{0,\nu-1} \bar r_{21}^\mu L_1^{\mu,\alpha-2} L_2^{\nu,\beta-2} e_1 e_2 - \sum_{\mu=1}^{\alpha-2} \sum_{\nu=1}^{\beta-1} \delta^{\mu\nu} L_1^{0,\mu-1} L_2^{0,\nu-1} \bar r^\mu L_1^{\mu,\alpha-2} L_2^{\nu,\beta-2} e_1 e_2 \\
            ={}& (r^0 - \delta_{1<\alpha<\beta} P + \tfrac{1}{2}\delta_{1<\alpha<\beta} + \tfrac{1}{2} \delta_{1<\beta<\alpha}) a_1^\alpha a_2^\beta
        \end{aligned}
    \end{equation*}
    where we used the identities
    \begin{equation*}
        \underline r^\alpha e_1 e_2 = 0, \qquad (\underline r^\alpha + \bar r^\alpha) e_2 = -\tfrac{1}{2} e_2, \qquad (\underline r^\alpha - \bar r_{21}^\alpha - P) e_1 = -\tfrac{1}{2} e_1.
    \end{equation*}
    The other brackets follow similarly, except that we have to use the additional identity
    \begin{equation*}
        L^0 \cdots L^{\alpha-2} \tilde L^{\alpha-1} \cdots \tilde L^{\ell-1} = \tilde L + \sum_{\rho=1}^{\alpha-1} a^\rho c^\rho,
    \end{equation*}
    which can be proven by a similar telescopic argument as in the proof of lemma \ref{lemma:hypLOp}. All in all, we obtain
    \begin{align}
        \{ a_1^\alpha,a_2^\beta \} ={}& (r^0 - \delta_{\alpha<\beta} P) a_1^\alpha a_2^\beta + (\delta_{1=\alpha<\beta} + \tfrac{1}{2} \delta_{\alpha>\beta>1} + \tfrac{1}{2}\delta_{1<\alpha<\beta}) a_1^\alpha a_2^\beta,
        \\[6pt]
        \{ a_1^\alpha,c_2^\beta \} ={}& {-c_2^\beta} \bar r^0 a_1^\alpha - (\tfrac{1}{2} \delta^{\alpha\beta} + \tfrac{1}{2} \delta_{1=\alpha < \beta} - \tfrac{1}{2} \delta_{\alpha > \beta=1}) a_1^\alpha c_2^\beta - \delta^{\alpha\beta} \sum_{\rho=1}^{\alpha-1} a_1^\rho c_2^\rho \\
        &+ \delta^{1\beta} Z a_1^\alpha \tilde L_2 - \delta^{\alpha\beta} Z e_1 \tilde L_2, \nonumber
        \\[6pt]
        \{ c_1^\alpha,c_2^\beta \} ={}& {-c_1^\alpha} c_2^\beta (\underline r^0 - \delta_{\alpha<\beta} P) - (\delta_{\alpha>\beta=1} + \tfrac{1}{2} \delta_{\alpha>\beta>1} + \tfrac{1}{2} \delta_{1<\alpha<\beta}) c_1^\alpha c_2^\beta \\
        &- \delta^{1\beta} c_1^\alpha Z \tilde L_2 + \delta^{1\alpha} c_2^\beta Z_{21} \tilde L_1, \nonumber
    \end{align}
    where $Z = \sum_{i=1}^N e_{ii} \otimes e_i^t$ and $P = \sum_{ij} e_{ij} \otimes e_{ji}$. Writing this in components, we arrive at the claimed brackets.
\end{proof}

\begin{remark}
    We again have the constraint $a_i^1 = 1$, which restricts the ``physical'' spin degrees of freedom of the hyperbolic spin RS model to $a^\alpha,c^\alpha$ for $\alpha > 1$ and notice that they form a quadratic Poisson algebra.
\end{remark}

\begin{corollary}
    \def\a{\alpha}
    \def\b{\beta}
    \def\sgn{\rm sgn}
    The rescaling $\hat a_i^\alpha \coloneq a_i^\alpha/\sum_\rho a_i^\rho, \hat c_i^\alpha \coloneq c_i^\alpha \sum_\rho a_i^\rho, \hat{ \tilde{L}}_{ij} \coloneq \frac{\sum_\rho \hat a_i^\rho \hat c_j^\rho}{Q_i^0/Q_j^\ell-1}$ reproduces the Poisson brackets from \cite{arutyunov:2019,fairon:2026}:
    \begin{align}
            \{ Q_i^0, \hat a_j^\alpha \} ={}& 0, \qquad \{ Q_i^0, \hat c_j^\alpha \} = \delta_{ij} Q_i^0 \hat c_j^\alpha, \nonumber \\[4pt]
            \{\hat{a}_i^{\a},\hat{a}_j^{\b} \}
            ={}&-(1-\delta_{ij})\frac{1}{2}\frac{Q_i^0+Q_j^0}{Q_i^0-Q_j^0}(\hat{a}_i^{\a}-\hat{a}_j^{\a})(\hat{a}_i^{\b}-\hat{a}_j^{\b})-\frac{\sgn(\beta-\a)}{2}\hat{a}_j^{\a}\hat{a}_i^{\b} \nonumber \\
            &-\sum_{\rho=1}^{\ell}\frac{\sgn(\a-\rho)}{2}\hat{a}_j^{\a}\hat{a}_j^{\b}\hat{a}_i^{\rho}+\sum_{\rho=1}^{\ell}\frac{\sgn(\b-\rho)}{2}\hat{a}_i^{\a}\hat{a}_i^{\b}\hat{a}_j^{\rho}
            -\sum_{\mu,\nu=1}^{\ell}\frac{\sgn(\mu-\nu)}{2}\hat{a}_i^{\mu}\hat{a}_j^{\nu} \hat{a}_i^{\a}\hat{a}_j^{\b}\, , \nonumber \\
            \{\hat{a}_i^{\a},\hat{c}_j^{\b} \}={}&\hat{a}_i^{\a} \hat{ \tilde{L}}_{ij}-(1-\delta_{ij})\frac{1}{2}\frac{Q_i^0+Q_j^0}{Q_i^0-Q_j^0}(\hat{a}_i^{\a}-\hat{a}_j^{\a})\hat{c}_j^{\b}+\hat{a}_i^{\a}\sum_{\rho=1}^{\beta-1}\hat{a}_i^{\rho}\hat{c}_j^{\rho}+\frac{1}{2}\hat{a}_i^{\a}\hat{a}_i^{\b}\hat{c}_j^{\b} \nonumber \\
            &-\delta_{\a\b}\bigg(\hat{ \tilde{L}}_{ij}+\frac{1}{2}\hat{a}_i^{\a}\hat{c}_j^{\b}+\sum_{\rho=1}^{\beta-1}\hat{a}_i^{\rho}\hat{c}_j^{\rho}\bigg)
            +\sum_{\rho=1}^{\ell}\frac{\sgn(\a-\rho)}{2}\hat{c}_{j}^{\b}\hat{a}_j^{\a}\hat{a}_i^{\rho}
            +\sum_{\mu,\nu=1}^{\ell}\frac{\sgn(\mu-\nu)}{2}\hat{a}_i^{\mu}\hat{a}_j^{\nu}\hat{a}_i^{\a}\hat{c}_j^{\b}\, , \displaybreak \nonumber \\
            \{\hat{c}_i^{\a},\hat{c}_j^{\b} \}={}&\frac{1}{2}(1-\delta_{ij})\frac{Q_i^0+Q_j^0}{Q_i^0-Q_j^0}(\hat{c}_i^{\a}\hat{c}_j^{\b}+\hat{c}_j^{\a}\hat{c}_i^{\b})-\hat{c}_i^{\a} \hat{ \tilde{L}}_{ij}+\hat{c}_j^{\b} \hat{ \tilde{L}}_{ji}+\frac{\sgn(\b-\a)}{2}\hat{c}_j^{\a}\hat{c}_i^{\b} \nonumber \\
            &-\hat{c}_i^{\a}\sum_{\rho=1}^{\b-1}\hat{a}_i^{\rho}\hat{c}_j^{\rho}+\hat{c}_j^{\b}\sum_{\rho=1}^{\a-1}\hat{a}_j^{\rho}\hat{c}_i^{\rho}
            -\sum_{\mu,\nu=1}^{\ell}\frac{\sgn(\mu-\nu)}{2}\hat{a}_i^{\mu}\hat{a}_j^{\nu} \hat{c}_i^{\a}\hat{c}_j^{\b}+\frac{1}{2}\hat{a}_j^{\a}\hat{c}_i^{\a}\hat{c}_j^{\b}-\frac{1}{2}\hat{a}_i^{\b}\hat{c}_{i}^{\a}\hat{c}_j^{\b}\, .
    \end{align}
\end{corollary}

\begin{remark}
    The fact that our Poisson brackets coincide with the ones in \cite{arutyunov:2019,fairon:2026} can be viewed as an instance of mirror symmetry. The principle of mirror symmetry posits a correspondence between the multiplicative quiver variety of the Jordan quiver considered in \cite{fairon:2026} and the $K$-theoretic Coulomb branch of the necklace quiver considered here.
\end{remark}

\subsection{Supplementary Poisson brackets}

In this section, we give supplementary Poisson brackets of certain variables introduced in \cite{arutyunov:1998}:
\begin{equation}
    S_i^{\alpha\beta} \coloneq c_i^\alpha a_i^\beta, \qquad f_{ij} \coloneq \sum_{\rho=1}^\ell a_i^\rho c_j^\rho.
\end{equation}
The variables $S_i^{\alpha\beta}$ were dubbed \emph{quantum current} in \cite{arutyunov:2025} due to the form of their quantum commutation relations, which are similar to the commutation relations of the quantum current defined in \cite{reshetikhin:1990}. The variables $f_{ij}$ are the \emph{collective spin variables}. They Poisson-commute with $Q_i^0$ according to
\begin{eqnarray}
    \{ Q_i^0,S_j^{\alpha\beta} \} = \delta_{ij} Q_i^0 S_j^{\alpha\beta}, \qquad \{ Q_i^0,f_{jk} \} = \delta_{ik} Q_i^0 f_{jk}.
\end{eqnarray}
We now present the Poisson brackets of the quantum current and collective spin variables of the hyperbolic spin RS model.

\begin{prop}
    \def\a{\alpha}
    \def\b{\beta}
    \def\sgn{\rm sgn}
    The hyperbolic quantum current satisfies the Poisson bracket
    \begin{equation}
        \begin{aligned}
            \{&S_i^{\a\b},S_j^{\mu\nu}\}=(1-\delta_{ij})\frac{1}{2}\frac{Q^0_i+Q^0_j}{Q^0_i-Q^0_j}(S_i^{\mu\beta}S_j^{\a\nu}+S_i^{\a\nu}S_j^{\mu\beta})
            +\frac{\sgn(\mu-\a)}{2}S_i^{\mu\b}S_j^{\a\nu}-\frac{\sgn(\nu-\b)}{2}S_i^{\a\nu}S_j^{\mu\b}
            \\
            &
            -\delta^{\b\mu}\bigg(\tfrac{1}{2}S_i^{\a\b}S_j^{\mu\nu}+\sum_{\rho=1}^{\b-1}S_i^{\a\rho}S_j^{\rho\nu}+\frac{\sum_{\rho=1}^{\ell}S_i^{\a\rho}S_j^{\rho\nu}}{Q_i^0/Q_j^{\ell}-1} \bigg) 
            +\delta^{\a\nu}\bigg(\tfrac{1}{2}S_i^{\a\b}S_j^{\mu\nu}+\sum_{\rho=1}^{\a-1}S_j^{\mu\rho}S_i^{\rho\beta}+\frac{\sum_{\rho=1}^{\ell}S_j^{\mu\rho}S_i^{\rho\beta}}{Q_j^0/Q_i^{\ell}-1} \bigg).
        \end{aligned}
    \end{equation}
\end{prop}

\begin{prop}
    \def\a{\alpha}
    \def\b{\beta}
    The collective spin variables satisfy the Poisson bracket
    \begin{eqnarray}
        \begin{aligned}
            \{ f_{ij}, f_{kl} \}
            &= (1-\delta_{ik})\bigg(\frac{(f_{kl}-f_{il})f_{ij}}{Q_i^0/Q_k^0-1}-\frac{(f_{kl}-f_{il})f_{kj}}{1-Q_k^0/Q_i^0}\bigg) + (1-\delta_{jl})\bigg(\frac{f_{il} f_{kj}}{Q_j^0/Q_l^0-1}+\frac{f_{ij}f_{kl}}{1-Q_l^0/Q_j^0}\bigg) \\
            &- (1-\delta_{kj})\frac{f_{ij}(f_{jl}-f_{kl})}{1-Q_j^0/Q_k^0}+(1-\delta_{il})\frac{f_{kl}(f_{lj}-f_{ij})}{1-Q_l^0/Q_i^0} \\
            &+ \frac{f_{ij} f_{il}}{Q_i^0/Q_l^\ell-1} -\frac{f_{ij} f_{jl}}{Q_j^0/Q_l^\ell-1} + \frac{f_{kl} f_{lj}}{Q_l^0/Q_j^\ell-1} - \frac{f_{kl} f_{kj}}{Q_k^0/Q_j^\ell-1} + \frac{f_{il} f_{kj}}{Q_k^0/Q_j^\ell-1} - \frac{f_{kj} f_{il}}{Q_i^0/Q_l^\ell-1}.
        \end{aligned}
    \end{eqnarray}
\end{prop}

\begin{remark}
    It can be checked that the rational degeneration of these Poisson brackets yields the Poisson bracket found in \cite{arutyunov:1998} and that the form of the Poisson brackets of the collective spin variables matches the conjecture put forth in \emph{loc.\ cit.}
\end{remark}

\section{Quantization of the trigonometric spin Ruijsenaars--Schneider model}

\section{The $K$-theoretic Coulomb branch algebra}\label{sec:CBAlg}

In this section, we construct generators of the abelianized quantized $K$-theoretic Coulomb branch algebra of the necklace quiver following \cite[remark 2.10]{tsymbaliuk:2023}. The generators can be expressed in terms of the exponentiated vacuum expectation values of abelianized gauge multiplet scalars $Q_i^\alpha$, where $\alpha \in \Z/\ell\Z$ labels a node of the quiver and $i = 1,\dots,N$ is a color index for the corresponding gauge group factor $U(N)_\alpha$, as well as the $q$-difference operators $P_i^\alpha$, which satisfy the commutation relations
\begin{equation}
    Q_i^\alpha Q_j^\beta = Q_j^\beta Q_i^\alpha, \qquad P_i^\alpha Q_j^\beta = q^{\delta_{ij} \delta^{\alpha\beta}} Q_j^\beta P_i^\alpha, \qquad P_i^\alpha P_j^\beta = P_j^\beta P_i^\alpha.
\end{equation}
The abelianized monopole operators of minuscule coweight/magnetic charge, which together with the exponentiated vacuum expectation values $Q_i^\alpha$ generate the abelianized quantized $K$-theoretic Coulomb branch algebra, can then be expressed in terms of the $q$-difference operators as
\begin{align}
    u_i^{\alpha+} &\coloneq \prod_{k=1}^N (Q_k^\alpha/Q_k^{\alpha+1})^{1/2} \frac{\prod_{k=1}^N (1-(\mu^\alpha)^{-1} Q_k^{\alpha+1}/Q_i^\alpha)}{\prod_{k(\neq i)} (1-Q_k^\alpha/Q_i^\alpha)} P_i^\alpha, \\
    u_i^{\alpha-} &\coloneq \prod_{k=1}^N (Q_k^\alpha/Q_k^{\alpha-1})^{1/2} \frac{\prod_{k=1}^N (1-\mu^{\alpha-1} Q_k^{\alpha-1}/Q_i^\alpha)}{\prod_{k(\neq i)} (1-Q_k^\alpha/Q_i^\alpha)} (P_i^\alpha)^{-1}.
\end{align}
We now describe their commutation relations. To this end, we define the structure function
\begin{equation}
    g_\kappa^\pm(z) \coloneq \frac{q^{\mp\kappa/2} z - 1}{z - q^{\mp\kappa/2}}.
\end{equation}
Given $(i,\alpha) \neq (j,\beta)$, we then find
\begin{align}
    u_i^{\alpha\pm} u_j^{\beta\pm} &= g_{\kappa^{\alpha\beta}}^\pm\big((q^{-1/2} \mu^\alpha)^{\delta^{\alpha+1,\beta}} (q^{-1/2} \mu^\beta)^{-\delta^{\alpha,\beta+1}} Q_{ij}^{\alpha\beta}\big) u_j^{\beta\pm} u_i^{\alpha\pm}, & (\ell > 2) \\
    u_i^{\alpha\pm} u_j^{\beta\pm} &= g_{2}^\pm\big(Q_{ij}^{\alpha\beta}\big) u_j^{\beta\pm} u_i^{\alpha\pm}, & (\alpha = \beta,\ell = 2) \\
    u_i^{\alpha\pm} u_j^{\beta\pm} &= g_{-1}^\pm\big(q^{-1/2}\mu^\alpha Q_{ij}^{\alpha\beta}\big) g_{-1}^\pm\big(q^{1/2}(\mu^\beta)^{-1} Q_{ij}^{\alpha\beta}\big) u_j^{\beta\pm} u_i^{\alpha\pm}, & (\alpha \neq \beta,\ell = 2) \\
    u_i^{0\pm} u_j^{0\pm} &= g_{2}^\pm\big(Q_{ij}^{00}\big) g_{-1}^\pm\big(q^{-1/2}\mu^0 Q_{ij}^{00}\big) g_{-1}^\pm\big(q^{1/2}(\mu^0)^{-1} Q_{ij}^{00}\big) u_j^{0\pm} u_i^{0\pm}, & (\ell = 1)
\end{align}
where we have used the abbreviation $Q_{ij}^{\alpha\beta} \coloneq Q_i^\alpha/Q_j^\beta$ and $\kappa^{\alpha\beta} \coloneq 2\delta^{\alpha\beta} - \delta^{\alpha+1,\beta} - \delta^{\alpha,\beta+1}$ is the Cartan matrix of the necklace quiver.

\section{The $L$-operator algebra}\label{sec:LOpAlgebra}

The main tool for obtaining the classical trigonometric spin RS model from the $K$-theoretic Coulomb branch algebra is the $L$-operator algebra introduced in \cite{arutyunov:2026}.
In this section 
we quantize this algebra. To every arrow $\alpha \to \alpha+1$ in the necklace quiver, we associate two $L$-operators $L^{\alpha\pm}$, which are $N \times N$ matrices with coefficients
\begin{align}
    L_{ij}^{\alpha+} &\coloneq \frac{q^{1/2}}{1-(\mu^\alpha)^{-1} Q_i^{\alpha+1}/Q_j^\alpha} u_j^{\alpha+}, \label{eq:LpOp} \\
    L_{ij}^{\alpha-} &\coloneq \frac{q^{-1/2}}{\mu^\alpha Q_i^\alpha/Q_j^{\alpha+1}-1} u_j^{\alpha+1,-}. \label{eq:LmOp}
\end{align}
In \cite{arutyunov:2026} it was shown that 
the Poisson brackets of the classical counterparts of these $L$-operators are quadratic in $L^{\alpha\pm}$ and can be written in terms of classical $r$-matrices. Assuming the operator ordering prescribed by equations \eqref{eq:LpOp} and \eqref{eq:LmOp}, we obtain the following commutation relations
\begin{align}
    R_{ab}^{\alpha+1,\beta+1}(q^{-1}) L_b^{\beta+} \bar R_{ab}^{\alpha+1,\beta}(q^{-1})^{-1} L_a^{\alpha+} &= L_a^{\alpha+} \bar R_{ba}^{\beta+1,\alpha}(q^{-1})^{-1} L_b^{\beta+} \underline R_{ab}^{\alpha\beta}(q^{-1}), \label{eq:oneSiteLpAlg} \\
    R_{ab}^{\alpha\beta}(q) L_b^{\beta-} \bar R_{ab}^{\beta+1,\alpha}(q)^{-1} L_a^{\alpha-} &= L_a^{\alpha-} \bar R_{ba}^{\beta,\alpha+1}(q)^{-1} L_b^{\beta-} \underline R_{ab}^{\alpha+1,\beta+1}(q). \label{eq:oneSiteLmAlg}
\end{align}
These relations provide a quantization of the aforementioned Poisson brackets.
Here, $a$ and $b$ denote auxiliary spaces and we introduce the following quantum $R$-matrices:
\begin{align}\nonumber
    R^{\alpha\beta}(q) \coloneq{}& 1 + \delta^{\alpha\beta} (1-q^{-1}) \times \\ & \Bigg[ \sum_{i \neq j} \bigg( \frac{1}{Q_i^\alpha/Q_j^\alpha-1} e_{ii} - \frac{1}{1-Q_j^\alpha/Q_i^\alpha} e_{ij} \bigg) \otimes (e_{jj}-e_{ji}) + X - X_{21} \Bigg], \\
    \bar R^{\alpha\beta}(q) \coloneq{}& 1 + \delta^{\alpha\beta} \bigg[ (q^{-1/2}-1) + \sum_{i \neq j} \frac{q^{1/2}-q^{-1/2}}{1 - q Q_j^\alpha/Q_i^\alpha} (e_{ii}-e_{ij}) \otimes e_{jj} + (q^{1/2}-q^{-1/2}) X \bigg], \\
    \underline R^{\alpha\beta}(q) \coloneq{}& 1 + \delta^{\alpha\beta} (1-q^{-1}) \sum_{i \neq j} \frac{1}{1-Q_j^\alpha/Q_i^\alpha} (e_{ij} \otimes e_{ji} - e_{ii} \otimes e_{jj}),
\end{align}
with $X = \sum_{i,j=1}^N e_{ij} \otimes e_{jj}$. These $R$-matrices are identical to the evaluation at zero of the spectral parameter dependent 
$R$-matrices found in \cite{arutyunov:2019b} for the spinless case up to a prefactor for $\bar R^{\alpha\beta}(q)$. It should be noted that $R^{\alpha\beta}(q)$ satisfies the Yang--Baxter equation
\begin{equation}
    R_{ab}^{\alpha\beta}(q) R_{ac}^{\alpha\gamma}(q) R_{bc}^{\beta\gamma}(q) = R_{bc}^{\beta\gamma}(q) R_{ac}^{\alpha\gamma}(q) R_{ab}^{\alpha\beta}(q),
\end{equation}
and $\underline R^{\alpha\beta}(q)$ satisfies the dynamical Yang--Baxter equation
\begin{equation}
    \underline R_{ab}^{\alpha\beta}(q) (P_b^\beta)^{-1} \underline R_{ac}^{\alpha\gamma}(q) P_b^\beta \underline R_{bc}^{\beta\gamma}(q) = (P_a^\alpha)^{-1} \underline R_{bc}^{\beta\gamma}(q) P_a^\alpha \underline R_{ac}^{\alpha\gamma}(q) (P_c^\gamma)^{-1} \underline R_{ab}^{\alpha\beta}(q) P_c^\gamma,
\end{equation}
where the role of the dynamical parameters is played by the variables $Q_i^\alpha$. Finally, $\bar R^{\alpha\beta}(q)$ is the corresponding dynamical twist, i.e.
\begin{equation}\label{eq:dynamicalTwistRel}
    \bar R^{\alpha\beta}(q) \underline R^{\alpha\beta}(q) = R^{\alpha\beta}(q) \bar R_{21}^{\alpha\beta}(q).
\end{equation}
These relations are shown diagrammatically in figure \ref{fig:Rmatrelations}.
\begin{figure}[t]
    \centering
    \begin{subfigure}{0.32\textwidth}
    \centering
    \begin{tikzpicture}[>=Stealth,thick,scale=0.75]
        \draw[->] (0.3,-1) -- (0.3,1);
        \draw[->] (-1,-1) -- (1,1);
        \draw[->] (1,-1) -- (-1,1);
        \node[below] at (-1,-1) {\footnotesize $\alpha$};
        \node[below] at (0.3,-1) {\footnotesize $\beta$};
        \node[below] at (1,-1) {\footnotesize $\gamma$};
        
        \node at (1.5,0) {$=$};
        
        \begin{scope}[xshift=3cm]
            \draw[->] (-0.3,-1) -- (-0.3,1);
            \draw[->] (-1,-1) -- (1,1);
            \draw[->] (1,-1) -- (-1,1);
            \node[below] at (-1,-1) {\footnotesize $\alpha$};
            \node[below] at (-0.3,-1) {\footnotesize $\beta$};
            \node[below] at (1,-1) {\footnotesize $\gamma$};
        \end{scope}
    \end{tikzpicture}
    \caption{}
    \end{subfigure}
    \begin{subfigure}{0.32\textwidth}
    \centering
    \begin{tikzpicture}[>=Stealth,thick,scale=0.75]
        \draw[gray!30!white,fill=gray!30!white] (-1,-1) rectangle (1,1);
        \draw[gray!30!white,fill=gray!30!white] (2,-1) rectangle (4,1);
    
        \draw[->] (0.3,-1) -- (0.3,1);
        \draw[->] (-1,-1) -- (1,1);
        \draw[->] (1,-1) -- (-1,1);
        \node[below] at (-1,-1) {\footnotesize $\alpha$};
        \node[below] at (0.3,-1) {\footnotesize $\beta$};
        \node[below] at (1,-1) {\footnotesize $\gamma$};
        \node[gray] at (0.6,0) {\scriptsize $Q$};
        
        \node at (1.5,0) {$=$};
        
        \begin{scope}[xshift=3cm]
            \draw[->] (-0.3,-1) -- (-0.3,1);
            \draw[->] (-1,-1) -- (1,1);
            \draw[->] (1,-1) -- (-1,1);
            \node[below] at (-1,-1) {\footnotesize $\alpha$};
            \node[below] at (-0.3,-1) {\footnotesize $\beta$};
            \node[below] at (1,-1) {\footnotesize $\gamma$};
            \node[gray] at (0.5,0) {\scriptsize $Q$};
        \end{scope}
    \end{tikzpicture}
    \caption{}
    \end{subfigure}
    \begin{subfigure}{0.32\textwidth}
    \centering
    \begin{tikzpicture}[>=Stealth,thick,scale=0.75]
        \draw[gray!30!white,fill=gray!30!white] (-1,-1) rectangle (1,0);
        \draw[gray!30!white,fill=gray!30!white] (2,-1) rectangle (4,0);
    
        \draw[->] (-0.7,-1) -- (1,0.7);
        \draw[->] (0.7,-1) -- (-1,0.7);
        \node[below] at (-0.7,-1) {\footnotesize $\alpha$};
        \node[below] at (0.7,-1) {\footnotesize $\beta$};
        \node[gray] at (0.5,-0.3) {\scriptsize $Q$};
        
        \node at (1.5,0) {$=$};

        \begin{scope}[xshift=3cm]
            \draw[->] (-1,-0.7) -- (0.7,1);
            \draw[->] (1,-0.7) -- (-0.7,1);
            \node[left] at (-1,-0.7) {\footnotesize $\alpha$};
            \node[right] at (1,-0.7) {\footnotesize $\beta$};
            \node[gray] at (0.82,-0.225) {\scriptsize $Q$};
        \end{scope}
    \end{tikzpicture}
    \caption{}
    \end{subfigure}
    \caption[Diagrammatic representation of the three relations satisfied by the $R$-matrices.]{Diagrammatic representation of the three relations satisfied by the $R$-matrices. (a) The Yang--Baxter equation for the $R$-matrix $R^{\alpha\beta}(q)$, which is depicted as a crossing of two lines labeled by $\alpha$ and $\beta$ on a white background. (b) The dynamical Yang--Baxter equation for the dynamical $R$-matrix $\underline R^{\alpha\beta}(q)$, which is depicted as a crossing of two lines labeled by $\alpha$ and $\beta$ on a gray background. The faces with a gray background are labeled by the dynamical parameters $Q$, with their values determined from the reference face (labeled $Q$) by $q$-shifting $Q^\alpha$ when passing over a line labeled $\alpha$. (c) The twist relation for the dynamical twist $\bar R^{\alpha\beta}(q)$, which is depicted as two lines labeled by $\alpha$ and $\beta$ crossing from a gray to a white background. Composition is read from bottom to top.}
    \label{fig:Rmatrelations}
\end{figure}

If we define the total $L$-operators to be
\begin{align}
    L^{\text{tot}+} \coloneq L^{\ell-1,+} \cdots L^{0,+}, \qquad L^{\text{tot}-} \coloneq L^{0-} \cdots L^{\ell-1,-},
\end{align}
the relations \eqref{eq:oneSiteLpAlg} and \eqref{eq:oneSiteLmAlg} together with equation \eqref{eq:dynamicalTwistRel} imply
\begin{equation}\label{eq:totalAlg}
    R_{ab}^{00}(q^{\mp}) L_b^{\text{tot}\pm} \bar R_{ab}^{00}(q^{\mp})^{-1} L_a^{\text{tot}\pm} = L_a^{\text{tot}\pm} \bar R_{ba}^{00}(q^{\mp})^{-1} L_b^{\text{tot}\pm} \underline R_{ab}^{00}(q^{\mp}),
\end{equation}
which exactly reproduces the $L$-operator algebra of the spinless model found in \cite{arutyunov:2019b}.

\section{Commuting Hamiltonians}\label{sec:Hamiltonians}

The commutation relations \eqref{eq:totalAlg} of the total $L$-operator coincide with those obtained in \cite{arutyunov:2019b} for the $L$-operator of the spinless model. It follows that a family of commuting Hamiltonians is given by the same trace formula as in \emph{loc.\ cit.}:
\begin{equation}\label{eq:traceFormula}
    H^\pm[n] \coloneq \operatorname{Tr}_{ab} P_{ab}^{t_b} L_a^{\text{tot}\pm} (\bar R_{ba}^{00}(q^{\mp})^{t_b} R_{ab}^{00}(q^{\mp})^{t_b} L_a^{\text{tot}\pm})^{n-1},
\end{equation}
where $P = \sum_{i,j=1}^N e_{ij} \otimes e_{ji}$ and the superscript $t_b$ denotes transposition in the auxiliary space $b$. In particular, we have $H^-[1] = \operatorname{Tr} L^{\text{tot}-}$, which coincides with the defining Hamiltonian of the trigonometric spin RS model introduced in \cite{arutyunov:2026} and generates at the classical level the equations of motion found by Krichever and Zabrodin \cite{krichever:1995}. We remark that in terms of the 4d $\mathcal{N}=2$ necklace quiver gauge theory, the Hamiltonian $H^-[1]$ corresponds to the 't Hooft line which has charge $\square$ under all gauge nodes \cite{maruyoshi:2021}.

\smallskip
The Hamiltonians \eqref{eq:traceFormula} can be rewritten in terms of the generalized Macdonald operators
\begin{equation}\label{eq:macdonald}
    \begin{aligned}
    S^+[n] &\coloneq t^{-n(n-1)/2} \!\!\!\!\!\!\!\!\!\! \sum_{\substack{I_0,\dots,I_{\ell-1} \subseteq \{ 1,\dots,N \} \\ |I_\alpha|=n}} \Bigg( \prod_{\alpha=0}^{\ell-1} \frac{\prod_{i \in I_{\alpha-1},k \notin I_\alpha} (Q_i^{\alpha-1}-(\mu^{\alpha-1})^{-1} Q_k^\alpha)}{\prod_{i \in I_{\alpha},k \notin I_\alpha} (Q_i^\alpha-Q_k^\alpha)} \Bigg) \prod_{\alpha=0}^{\ell-1} \prod_{i \in I_\alpha} P_i^\alpha, \\
    S^-[n] &\coloneq t^{n(n-1)/2} \!\!\!\!\!\! \sum_{\substack{I_0,\dots,I_{\ell-1} \subseteq \{ 1,\dots,N \} \\ |I_\alpha|=n}} \Bigg( \prod_{\alpha=0}^{\ell-1} \frac{\prod_{i \in I_{\alpha+1},k \notin I_\alpha} (Q_i^{\alpha+1}-\mu^\alpha Q_k^\alpha)}{\prod_{i \in I_{\alpha},k \notin I_\alpha} (Q_i^\alpha-Q_k^\alpha)} \Bigg) \prod_{\alpha=0}^{\ell-1} \prod_{i \in I_\alpha} (P_i^\alpha)^{-1}.
    \end{aligned}
\end{equation}
Specifically, we have the identification
\begin{equation}\label{eq:determinant}
    H^\pm[n] = q^{\pm 1/2}
    \begin{vmatrix}
        S^\pm[1] & 1 & 0 & \cdots & 0 \\
        [2]_{q^{\pm 1}} S^\pm[2] & S^\pm[1] & 1 & \cdots & 0 \\
        \vdots & \vdots & \vdots & \ddots & \vdots \\
        [n]_{q^{\pm 1}} S^\pm[n] ~&~ S^\pm[n-1] ~&~ S^\pm[n-2] & \cdots & S^\pm[1]
    \end{vmatrix},
\end{equation}
where we make use of the $q$-numbers $[n]_q = \sum_{k=0}^{n-1} q^k$. This suggests that $S^\pm[n]$ are the elementary symmetric counterparts of the power sums $H^\pm[n]$. Indeed, we find
\begin{equation}\label{eq:spectralCurve}
    : {\det(q^{\mp\ell/2} L^{\text{tot}\pm}+z)} : \, = \sum_{n=0}^N z^{N-n} S^\pm[n],
\end{equation}
where $::$ means that the $P_i^\alpha$ are normal ordered to the right of the $Q_i^\alpha$. Note that introducing a spectral parameter dependence in $L^{\text{tot}\pm}$ would turn equation \eqref{eq:spectralCurve} into the quantum spectral curve of the trigonometric spin RS model.

To finish the discussion of the Hamiltonians, let us note that we can build generators of the 1+1d Poincaré Lie algebra from the lowest Hamiltonians $H^\pm[1]$ as follows:
\begin{equation}
    \mathbf{H} \coloneq H^+[1] + H^-[1], \qquad \mathbf{P} \coloneq H^+[1] - H^-[1], \qquad \mathbf{B} \coloneq \sum_{i=1}^N \log Q_i^0.
\end{equation}
They obey the expected relations
\begin{equation}
    [\mathbf{H},\mathbf{P}] = 0, \qquad [\mathbf{H},\mathbf{B}] = \log(q) \mathbf{P}, \qquad [\mathbf{P},\mathbf{B}] = \log(q) \mathbf{H},
\end{equation}
and the boost generator $\mathbf{B}$ more generally satisfies
\begin{equation}
    [H^\pm[n],\mathbf{B}] = \pm n \log(q) H^\pm[n], \qquad [S^\pm[n],\mathbf{B}] = \pm n \log(q) S^\pm[n].
\end{equation}
This repeats the known pattern of the spinless trigonometric RS model.

\section{Quantum loop algebra}\label{sec:quantumLoopAlg}

In this section, we exhibit a quantum loop algebra structure inside the $K$-theoretic Coulomb branch algebra. As discussed in the introduction, the general theory of \cite[appendix B]{braverman:2018b} and the technology developed in \cite{tsymbaliuk:2023} suggest that the $K$-theoretic Coulomb branch of the necklace quiver should be identified with the $N$-truncated quantum toroidal algebra $U_{q,t}^{(N)}(\ddot{\mathfrak{gl}}_\ell)$. Specializing to zero central charge, the quantum toroidal algebra has two subalgebras isomorphic to the $N$-truncated quantum loop algebra $U_q^{(N)}(\dot{\mathfrak{gl}}_\ell)$, called \emph{horizontal} and \emph{vertical}. Here, we exhibit the horizontal quantum loop subalgebra.

\smallskip
The quantum loop algebra $U_q(\dot{\mathfrak{gl}}_\ell)$ has various useful presentations. Firstly, it may be presented in the Drinfeld--Jimbo generators
\begin{equation}
    E^\alpha, \qquad (t^\alpha)^{\pm 1}, \qquad F^\alpha,
\end{equation}
which satisfy relations analogous to the relations for $U_q(\mathfrak{gl}_\ell)$ \cite{finkelberg:2019}:
\begin{align}
    (t^\alpha)^{\pm 1} (t^\alpha)^{\mp 1} &= 1, \label{eq:DJrel1} \\[5pt]
    [t^\alpha,t^\beta] &= 0, \\
    [E^\alpha,F^\beta] &= \delta^{\alpha\beta} \frac{K^\alpha-(K^\alpha)^{-1}}{q^{1/2}-q^{-1/2}}, \\
    E^\alpha t^\beta &= q^{(-\delta^{\alpha\beta}+\delta^{\alpha+1,\beta})/2} t^\beta E^\alpha \\[5pt]
    F^\alpha t^\beta &= q^{(\delta^{\alpha\beta}-\delta^{\alpha+1,\beta})/2} t^\beta F^\alpha, \\[5pt]
    E^\alpha E^\beta &= E^\beta E^\alpha, \qquad \alpha-\beta \neq 0,\pm1, \\[5pt]
    F^\alpha F^\beta &= F^\beta F^\alpha, \qquad \alpha-\beta \neq 0,\pm1,
\end{align}
where $K^\alpha \coloneq t^\alpha (t^{\alpha-1})^{-1}$, as well as the $q$-Serre relations
\begin{align}
    (E^\alpha)^2 E^{\alpha\pm1} + E^{\alpha\pm1} (E^\alpha)^2 &= (q^{1/2}+q^{-1/2}) E^\alpha E^{\alpha\pm1} E^\alpha, \\[5pt]
    (F^\alpha)^2 F^{\alpha\pm1} + F^{\alpha\pm1} (F^\alpha)^2 &= (q^{1/2}+q^{-1/2}) F^\alpha F^{\alpha\pm1} F^\alpha,
\end{align}
in the case $\ell > 2$, and
\begin{align}
    (E^\alpha)^3 E^{\alpha\pm1} - E^{\alpha\pm1} (E^\alpha)^3 &= (q^{-1} + 1 + q) \Big[ (E^\alpha)^2 E^{\alpha\pm1} E^\alpha - E^\alpha E^{\alpha\pm1} (E^\alpha)^2 \Big], \\[5pt]
    (F^\alpha)^3 F^{\alpha\pm1} - F^{\alpha\pm1} (F^\alpha)^3 &= (q^{-1} + 1 + q) \Big[ (F^\alpha)^2 F^{\alpha\pm1} F^\alpha - F^\alpha F^{\alpha\pm1} (F^\alpha)^2 \Big], \label{eq:DJrel2}
\end{align}
in the case $\ell = 2$. The case $\ell=1$ does not have Serre relations.

\smallskip
Secondly, the RTT presentation \cite{molev} of the quantum loop algebra $U_q(\dot{\mathfrak{gl}}_\ell)$ is given in terms of generator matrices
\begin{equation}
    T^\pm(u) = \sum_{\alpha,\beta=1}^\ell \sum_{n \geq 0} T^\pm[n]^{\alpha\beta} \otimes u^{\mp n} e^{\alpha\beta} \in U_q(\dot{\mathfrak{gl}}_\ell) \otimes \operatorname{End}(\C^\ell)\llbracket u^{\mp 1} \rrbracket,
\end{equation}
which satisfy the additional requirements 
\begin{equation}
    T^\pm[0]^{\alpha\alpha} T^\mp[0]^{\alpha\alpha} = 1, \qquad T^\pm[0]^{\alpha\beta} = 0, \qquad \alpha \gtrless \beta,
\end{equation}
as well as the RTT relations
\begin{align}\label{eq:RTT}
    R^{ab}(u/v) T^\pm(u)^a T^\pm(v)^b &= T^\pm(v)^b T^\pm(u)^a R^{ab}(u/v), \\
    R^{ab}(u/v) T^-(u)^a T^+(v)^b &= T^+(v)^b T^-(u)^a R^{ab}(u/v),
\end{align}
where we have used the standard unitary trigonometric $R$-matrix for $\mathfrak{gl}_\ell$:
\begin{equation}\label{eq:trigRMat}
    \begin{aligned}
        R(u) ={}& \frac{q^{1/2}-q^{-1/2} u}{q^{-1/2}-q^{1/2} u} \sum_\alpha e^{\alpha\alpha} \otimes e^{\alpha\alpha} + \frac{1-u}{q^{-1/2}-q^{1/2} u} \sum_{\alpha\neq\beta} e^{\alpha\alpha} \otimes e^{\beta\beta} \\
        &+ \frac{q^{1/2}-q^{-1/2}}{q^{-1/2}-q^{1/2} u} \sum_{\alpha>\beta} e^{\alpha\beta} \otimes e^{\beta\alpha} + \frac{q^{1/2} u-q^{-1/2} u}{q^{-1/2}-q^{1/2} u} \sum_{\alpha<\beta} e^{\alpha\beta} \otimes e^{\beta\alpha}.
    \end{aligned}
\end{equation}
The isomorphism between the Drinfeld--Jimbo and the RTT presentation is explicitly given in \cite{finkelberg:2019} for the case of $U_q(\mathfrak{gl}_\ell)$ and has a natural generalization for the case of the quantum loop algebra $U_q(\dot{\mathfrak{gl}}_\ell)$.

\smallskip
We now build a representation of the quantum loop algebra $U_q(\dot{\mathfrak{gl}}_\ell)$ from the abelianized monopole operators $u_i^\alpha$. Indeed, it may be checked that the generators
\begin{align}
    t^\alpha &\coloneq \prod_{k=1}^N ((\mu^{\alpha-1})^{-1} Q_k^\alpha/Q_k^{\alpha-1})^{1/2}, \label{eq:DJgen1} \\
    E^{\alpha} &\coloneq \frac{(\mu^{\alpha-1})^{-N/2}}{q^{-1/2}-q^{1/2}} \sum_{i=1}^N u_i^{\alpha+}, \\
    F^{\alpha} &\coloneq \frac{(\mu^\alpha)^{N/2}}{q^{1/2}-q^{-1/2}} \sum_{i=1}^N u_i^{\alpha-}, \label{eq:DJgen2}
\end{align}
satisfy the relations \eqref{eq:DJrel1}--\eqref{eq:DJrel2}. More generally, we can use the $L$-operator algebra to define the combinations
\begin{align}
    E^{\alpha,\beta} \coloneq (\mu^{\alpha-1})^{-1} u^{\beta+} L^{\beta-1,+} \cdots L^{\alpha+} e, \qquad F^{\alpha,\beta} \coloneq \mu^\beta u^{\alpha-} L^{\alpha-} \cdots L^{\beta-1,-} e,
\end{align}
where we use the covector $u^{\alpha\pm} \coloneq (u_1^{\alpha\pm},\dots,u_N^{\alpha\pm})$ and the vector $e = (1,\dots,1)^t$. In terms of the 4d $\mathcal{N}=2$ necklace quiver gauge theory, these may be thought of as monopole operators with fundamental charge under the gauge nodes $\alpha$ to $\beta$, see \cite[\S 6.6.1]{bullimore:2015}. From the integrable spin chain perspective, they may be thought of as the propagation of an excitation along the spin chain from site $\alpha$ to site $\beta$. Using the isomorphism between the Drinfeld--Jimbo and the RTT presentation \cite{finkelberg:2019}, we then find the zeroth and first modes of the generator matrices $T^\pm(u)$ to be of the form
\begin{equation}
    T^+[0]^{\alpha\beta} =
    \begin{cases}
        (t^\alpha)^{-1} E^{\alpha,\beta-1}, & \alpha < \beta,  \\
        (t^\alpha)^{-1}, & \alpha = \beta, \\
        0, & \alpha > \beta,
    \end{cases} \qquad \qquad 
    T^-[0]^{\alpha\beta} =
    \begin{cases}
        F^{\beta,\alpha-1} t^\beta, & \alpha > \beta, \\
        t^\beta, & \alpha = \beta, \\
        0, & \alpha < \beta,
    \end{cases}
\end{equation}
and
\begin{align}
    T^+[1]^{\alpha\beta} &= q^{-(1-\delta^{\alpha\beta})/2} (t^\alpha)^{-1} E^{\alpha,\beta+\ell-1}, \qquad & T^-[1]^{\alpha\beta} &= q^{(1-\delta^{\alpha\beta})/2} F^{\beta,\alpha+\ell-1} t^\beta.
\end{align}
One may check that the RTT relations \eqref{eq:RTT} of the quantum loop algebra are indeed satisfied to first order in $u$ and $v$. Now, we can use the quantum Leibniz formula
\begin{equation}
    \operatorname{qdet}^\pm(u) = \sum_{n=0}^\infty \operatorname{qdet}^\pm[n] u^{\mp n} = \sum_{\sigma \in S_\ell} (-q^{1/2})^{\ell(\sigma)} T^{\sigma(1)1\pm}(u) \cdots T^{\sigma(\ell)\ell\pm}(q^{-\ell+1} u),
\end{equation}
with $S_\ell$ the symmetric group on $\ell$ objects, to extract the first modes of the quantum determinant and find
\begin{align}\label{eq:qdet}
    \operatorname{qdet}^\pm[0] = t^{\pm N/2}, \qquad \operatorname{qdet}^\pm[1] = q^{\mp 1/2} (1-t^{\mp 1}) H^\pm[1].
\end{align}
It follows that the Bethe subalgebra of the horizontal quantum loop algebra $U_q(\dot{\mathfrak{gl}}_\ell)$ (which includes the quantum determinant) provides a large algebra of commuting Hamiltonians including the defining Hamiltonian $H^-[1]$. By a naive counting, the image of the Bethe subalgebra in the $N$-truncated quotient algebra $U_q^{(N)}(\dot{\mathfrak{gl}}_\ell)$ has $N\ell$ algebraically independent generators, which would provide enough commuting Hamiltonians for Liouville integrability\footnote{We thank J. Lamers and M. Vasilev for this insight.}, since the $K$-theoretic Coulomb branch algebra is generated by the $2N\ell$ algebraically independent generators $Q_i^\alpha$ and $P_i^\alpha$.

\smallskip For the first few generalized Macdonald operators \eqref{eq:macdonald}, one can verify that they commute with the Drinfeld--Jimbo generators \eqref{eq:DJgen1}--\eqref{eq:DJgen2}. We therefore conjecture that the Hamiltonians $H^\pm[n]$ are central in the horizontal quantum loop algebra $U_q(\dot{\mathfrak{gl}}_\ell)$ and hence lie in the horizontal Heisenberg subalgebra defined in \cite{wen:2025}. The theory of \emph{loc.\ cit.}\ would suggest that eigenstates of the Hamiltonians can be given by wreath Macdonald polynomials.

\section{Quantization of physical spin variables}\label{sec:spinVars}

Following the formulas for  the classical case elaborated in \cite{arutyunov:2026}, the spin variables used by Krichever and Zabrodin \cite{krichever:1995} are defined by the expressions
\begin{equation}
    a^\alpha \coloneq L^{0-} \cdots L^{\alpha-1,-} e, \qquad c^\alpha \coloneq u^{\alpha+1,-} \tilde L^{\alpha+1,-} \cdots \tilde L^{\ell-1,-}
\end{equation}
for $\alpha = 0,\dots,\ell-1$, where we use the modified $L$-operator $\tilde L^{\alpha-} \coloneq \mu^\alpha Q^\alpha L^\alpha (Q^{\alpha+1})^{-1}$, which satisfies the equation\footnote{This can be seen as the trigonometric analog of the quantum moment map equations of the chainsaw quiver variety \cite[equation (34)]{finkelberg:2014}.}
\begin{equation}
    q^{1/2} L^{\alpha-} - q^{-1/2} \tilde L^{\alpha-} = e u^{\alpha+1,-}.
\end{equation}
The same telescopic argument as in \cite{arutyunov:2026} yields a similar identity for the total $L$-operators
\begin{equation}\label{eq:Ltot}
        \sum_{\alpha \in \Z/\ell\Z} q^\alpha a^\alpha c^\alpha
        = q^{\ell-1/2} L^{\text{tot},-} - q^{-1/2} \tilde L^{\text{tot},-},
\end{equation}
with $\tilde L^{\text{tot}-} \coloneq \tilde L^{0,-} \cdots \tilde L^{\ell-1,-} = t Q^0 L^{\text{tot}-} (Q^0)^{-1}$. We can solve equation \eqref{eq:Ltot} for the matrix coefficients, obtaining
\begin{equation}
    L_{ij}^{\text{tot}-} = \frac{q^{1/2-\ell}}{1- tq^{1-\ell} Q_i^0/Q_j^0} \sum_{\alpha \in \Z/\ell\Z} q^\alpha a_i^\alpha c_j^\alpha, \qquad \tilde L_{ij}^{\text{tot}-} = \frac{q^{1/2}}{t^{-1} q^{\ell-1} Q_j^0/Q_i^0-1} \sum_{\alpha \in \Z/\ell\Z} q^\alpha a_i^\alpha c_j^\alpha.
\end{equation}

We now comment on the commutation relations of the \emph{physical} spin variables $a^\alpha$ and $c^\alpha$, meaning we restrict ourselves to spin states with index $\alpha \neq 0$. We obtain:
\begin{align}
    Q_i^0 a_j^\alpha &= a_j^\alpha Q_i^0, \label{eq:spinRel1} \\[6pt]
    Q_i^0 c_j^\alpha &= q^{\delta_{ij}} c_j^\alpha Q_i^0, \\[6pt]
    a_a^\alpha a_b^\alpha &= R_{ab}^{00}(q) a_b^\alpha a_a^\alpha, \\[6pt]
    a_a^\alpha a_b^\beta &= q^{1/2} R_{ab}^{00}(q) a_b^\beta a_a^\alpha, & (\alpha < \beta) \\[-2pt]
    c_a^\alpha \bar R_{ba}^{00}(q)^{-1} a_b^\alpha &= q^{-1/2} a_b^\alpha c_a^\alpha - q^{-\alpha} (1-q^{-1}) \bigg[ \Delta_{ab} \tilde L_a^{\text{tot}-} + q^{1/2} \sum_{\rho=0}^{\alpha-1} q^\rho a_b^\rho c_a^\rho \bigg], \\[-3pt]
    c_a^\alpha \bar R_{ba}^{00}(q)^{-1} a_b^\beta &= a_b^\beta c_a^\alpha, & (\alpha < \beta) \\[6pt]
    c_a^\alpha c_b^\alpha &= c_b^\alpha c_a^\alpha \underline R_{ab}^{00}(q)^{-1}, \\[6pt]
    c_a^\alpha c_b^\beta &= q^{-1/2} c_b^\beta c_a^\alpha \underline R_{ab}^{00}(q)^{-1}, & (\alpha < \beta) \label{eq:spinRel2}
\end{align}
with $\Delta = \sum_{i=1}^N e_i^t \otimes e_i$.

\smallskip
The commutation relations \eqref{eq:spinRel1}--\eqref{eq:spinRel2} show how the spin variables commute given definite spin states $\alpha,\beta$. Physically, however, it makes more intuitive sense to ask how the spin variables of two given particles $i,j$ interchange. This question indeed has a nice answer in the case of the physical spin vectors $c_i$, where we observe a simple exchange relation:
\begin{equation}\label{eq:spinExchange}
    c_i^a c_i^b = R^{ab}(-1) c_i^b c_i^a, \qquad c_i^a c_j^b = R^{ab}(Q_j^0/Q_i^0) c_j^b c_i^a, \qquad (i \neq j)
\end{equation}
with $R(u)$ the standard trigonometric $R$-matrix \eqref{eq:trigRMat} for $\mathfrak{gl}_{\ell-1}$. However, a similar formula for the physical spin covectors $a_i$ is not available, repeating the pattern remarked upon after proposition 3.6 of \cite{arutyunov:2025} in the rational degeneration.

\smallskip
We now finish the discussion of the quantization of the physical spin variables by giving their quantum equations of motion (the Heisenberg equations). This can be done simply by utilizing the commutation relations \eqref{eq:oneSiteLmAlg}. For example, we can compute
\begin{align}
    a^\alpha H^-[1] = \operatorname{Tr}_a a_b^\alpha L_a^{\text{tot}-} = q^{1/2} \operatorname{Tr}_a R_{ba}^{00}(q) L_a^{\text{tot}-} \bar R_{ba}^{00}(q)^{-1} a_b^\alpha,
\end{align}
where on the way we have used the identity $\underline R_{ba}^{\beta\alpha}(q)^{-1} e_b = q^{1/2} \bar R_{ab}^{\alpha\beta}(q)^{-1} e_b$. Writing these identities in components, the Heisenberg equations become
\begin{align}
    \dot{Q}_i^0=   [H^-[1],Q_i^0] &= (q^{-1}-1) Q_i^0 L_{ii}^{\text{tot}-}, \label{eq:eoms1} \\[8pt]
    \dot{a}_i^\alpha=  [H^-[1],a_i^\alpha] &= (q^{-1}-1) \sum_{j(\neq i)} L_{ij}^{\text{tot}-} (a_j^\alpha-a_i^\alpha) \frac{1}{1-Q_j^0/Q_i^0}, \\[-2pt]
    \dot{c}_i^\alpha=  [H^-[1],c_i^\alpha] &= (q^{-1}-1) \sum_{j(\neq i)} \bigg( \frac{1}{1-Q_j^0/Q_i^0} c_i^\alpha L_{ij}^{\text{tot}-} - \frac{1}{1-Q_i^0/Q_j^0} c_j^\alpha L_{ji}^{\text{tot}-} \bigg), \label{eq:eoms2}
\end{align}
which give the natural operator ordering for the operator form of the classical equations of motion by Krichever and Zabrodin \cite{krichever:1995}.


\subsection{Schur quantization}

\section{Elliptic spin Ruijsenaars--Schneider model from 5d $\mathcal{N}=1$ gauge theory}

\cite{murinov:2025}

\section{Class $\mathcal{S}_k$ and the elliptic spin Ruijsenaars--Schneider model}

\section{Brane construction}

\cite{maruyoshi:2021}